% ============================================================
%  Georg Brandes, Dualismen i vor nyeste Philosophie
%  Kjøbenhavn: Forlagt af den Gyldendalske Boghandel (F. Hegel).
%  I. Cohens Bogtrykkeri.  1866.   (76 pp., antiqua)
%  TRANSCRIPTION (Danish)
%
%  ---- SOURCES -----------------------------------------------------------
%  This edition was made from TWO scans of two different physical copies,
%  and uses both.  Identity is the sha256, not the filename.
%
%  COPY-TEXT / image of record, and witness 1:
%      bibliotek/Brandes, Georg/Dualismen_i_vor_nyeste_philosophie.pdf
%      Google Books (British Museum copy), 82 pp, 600 dpi bitonal JBIG2,
%      with an embedded text layer.
%      sha256 62393babbe7a6d8f5c198ba5ba60716f1ca278f9c4f154b031695a661d2de63a
%
%  Witness 2 (never copy-text; an independent machine reading):
%      bibliotek/Brandes, Georg/dualismen-kb.pdf
%      Royal Danish Library, 89 pp, 250 ppi colour, ABBYY Recognition Server.
%      sha256 21c194347c101781b33a4181981fd29a2ae8f7b0938168f982c83f67d14e1ad2
%
%  Why two.  The copy-text scan is bitonal JBIG2, and JBIG2's symbol-dictionary
%  mode is the one compression scheme known to substitute whole glyphs rather
%  than degrade them.  A substitution invented by a compressor cannot also
%  appear in a 250 ppi JPEG of a different copy, so anything the two layers
%  agree on was not manufactured between them.  Every batch of this book was
%  diffed against BOTH layers before it was spliced.  The two-witness set-up
%  was tested before it was trusted: nine disagreements between the layers were
%  adjudicated at 350--900 dpi during set-up, and in every one the 600 dpi image
%  was clean and correct, the error lying in one OCR engine.  No JBIG2 glyph
%  substitution was found in this scan.
%
%  ---- PAGE MAP ----------------------------------------------------------
%  see pagemap.py — both offsets are uniform end to end.
%      printed 1--76  =  PDF +  4   (Google, PDF  5--80)
%      printed 1--76  =  PDF +  7   (KB,     PDF  8--83)
%
%  Verified 2026-09-20 AT THE IMAGE, not from a text layer: the folio was read
%  off the Google scan on printed 8, 9, 15, 16, 17, 24, 26, 39, 40, 41, 46, 48,
%  54, 56, 59, 60, 61, 69, 70, 71, 74, 75, 76 and off the KB scan on 8, 15, 16,
%  39, 40, 41, 75, 76; every numeral matched.  Corroborated independently by
%  aligning the two text layers at +4/+7: mean per-page token agreement 0.966
%  (min 0.831) across all 70 body pages, which a one-leaf drift anywhere would
%  collapse to near zero.  No rescanned leaf, no plate.
%
%  COLLATION.  The arabic sequence counts the preliminary leaves, but none of
%  them carries a printed folio, and the four chapter openings drop theirs too.
%      printed  1  PDF  5   half-title
%      printed  2  PDF  6   blank (British Museum stamp)
%      printed  3  PDF  7   title page
%      printed  4  PDF  8   blank (British Museum stamp)
%      printed  5  PDF  9   Indhold
%      printed  6  PDF 10   blank
%      printed  7  PDF 11   body begins, chapter 1
%      printed 76  PDF 80   last page of text, ending with the Efterskrift
%      PDF 81--82           publisher's advertisement list, outside the sequence
%  The first folio actually printed is 8, on PDF 12.  There is no dedication,
%  motto, epigraph, preface or errata leaf: the preliminaries are exactly the
%  three printed pages above.  Printer's signatures 2--6 fall at pp. 17, 33,
%  49, 65.
%
%  ---- EDITORIAL STANCE --------------------------------------------------
%  Diplomatic.  Nineteenth-century orthography exactly as printed
%  (Philosophie, høiere, Videnskaben, capitalised nouns, aa).
%
%  THE BOOK IS INCONSISTENT WITH ITSELF and is not normalised in either
%  direction.  Audited at the image, occurrence by occurrence:
%      Deel  p. 70          against  Del   pp. 19, 27, 38, 45, 60, 66, 74
%      deels p. 59          against  dels  p. 42
%      heelt p. 61          against  helt  pp. 27, 40, 56, 58
%      naturligviis p. 59 ONLY; the other eleven occurrences read naturligvis
%  Settled archaic forms throughout: een, eet, seer, seet, veed, Eenhed, puur.
%  Forms that look archaic but are not: kvantvis, ligervis, kvæles.
%  Steen / Been / Meer do not occur in this book at all.
%  Printer's errors are transcribed AS PRINTED and logged in a % comment at
%  the site, never silently corrected — so the quote-balance count may drift
%  off zero on purpose; each such drift is logged below and in RESUME-NOTES.
%
%  EMPHASIS.  Letterspacing (Sperrsatz) is this book's ONLY emphasis device;
%  there is no italic anywhere in it, and Brandes letterspaces whole clauses,
%  not just single words.  Letterspacing is rendered \emph{}.  NEITHER TEXT
%  LAYER RECORDS IT — both flatten it — so every \emph{} in this file was
%  decided at the image, assisted by spacing.py's mechanical measurement.
%
%  FOOTNOTES.  Marked *) and, where a page carries two, **).  The mark is
%  therefore NOT constant in this book, unlike Evangelietroen og Theologien,
%  and \thefootnote is redefined per note at the site (see below).  The note
%  on printed p. 11 runs over onto p. 12.
%
%  FOREIGN MATTER.  French (trinité chrétienne p. 56, n'observe p. 65) and
%  Pièce (p. 76) are set in the ordinary roman, with NO contrasting face, so
%  they take no \textit{}.  There is no Greek in this book.
%
%  ---- HOW THIS BOOK'S ACCURACY WAS ESTABLISHED --------------------------
%  Word-accuracy: every batch diffed against BOTH witnesses before splicing
%  (ocrdiff.py).  7 batches, 62 candidates (58 Google, 4 KB), 0 transcription
%  errors corrected, 62 the witness's fault, 0 unresolved.  No divergence
%  appeared in both diffs at once.  See TRANSCRIPTION-PLAYBOOK.md §3 step 4.
%
%  THE INSTRUMENT WAS TESTED BEFORE IT WAS TRUSTED.  A clean diff means
%  nothing unless the diff can fail, so three faults of the classes that
%  matter were planted in the finished pp. 7--15 fragment and both diffs
%  re-run.  Result: a dropped „ikke“ was caught by BOTH witnesses, and a
%  substituted content word (Videnskaben -> Videnskabet) by both.  A
%  MODERNISED ARCHAIC FORM (naturligviis -> naturligvis) WAS CAUGHT BY
%  NEITHER, because both layers modernise that word themselves.  That is the
%  method's real blind spot: it is blind exactly where both machine readings
%  share an error, and no amount of diffing will show it.
%
%  So the blind spot was measured directly rather than assumed small.  A
%  dedicated orthographic audit read every candidate at the image and found:
%    * the doubled-vowel family is MIXED in this book and both layers report
%      it correctly (the table under EDITORIAL STANCE above);
%    * naturligviis is printed on p. 59 and on p. 59 only, of twelve
%      occurrences -- the Google layer is wrong there and right elsewhere;
%    * the class both layers DO share is capital Æ, which Google reads as E
%      or A at pp. 19, 29, 36, 64, 70, 73, and æ in Propædeutik in the notes
%      to pp. 55 and 59.  Every one was read at the image and is correct here.
%    * neither layer is reliable on punctuation anywhere in this book, so
%      every dash, quotation mark and stop in this file came from the image.
%  The two layers agree on 91.1% of word slots; they can be trusted on the
%  book's letters nearly everywhere and on its punctuation nowhere.
%
%  INDEPENDENT COLLATION.  Because every batch reported 0 corrections -- the
%  pattern the playbook says to distrust -- the result was measured instead of
%  accepted.  14 of the 70 body pages (20%) were drawn by a SEEDED random
%  generator (seed 18661220, so the sample is reproducible and was not chosen
%  after the fact) and read at the image, word by word, against this file by
%  two agents that had not transcribed the book and were told to assume it
%  wrong: pp. 18, 31, 32, 33, 38, 47, 54, 57, 58, 61, 63, 64, 67, 74.
%      Result: 0 discrepancies.  0 word-level, 0 orthographic, 0 punctuation,
%      0 emphasis, 0 structural.  Measured rate 0.00 errors per page; with 14
%      clean pages the 95% upper bound is ~0.21 errors per page.
%  A third check settled seven doubled-word sites found by a mechanical sweep
%  (all seven are Brandes's own repetition or carry punctuation between them;
%  all correct as transcribed) and re-verified emphasis boundaries at p. 26,
%  p. 76 and two randomly drawn sites.
%
%  PAGE JOINTS.  splice.py separates fragments with a blank line, which TeX
%  reads as a paragraph break, so a sentence running across a batch boundary
%  silently acquires a break the book does not have.  joints.py found three
%  (pp. 16, 36, 66); each was repaired and then VERIFIED AT THE IMAGE, since
%  the repair itself was made from the file and not from the page.  All three
%  pages open flush at the body margin -- measured at 200 dpi as +1, 0 and 0
%  px against a paragraph indent of +56 to +58 px elsewhere in the book -- so
%  the text does run on and the repairs are right.  No word is broken across
%  those three turns.  Intra-page paragraphing at each joint matches the image
%  and nothing was swallowed.  joints.py still reports pp. 3 and 5: those are
%  the front-matter \begin{center} blocks and are FALSE POSITIVES -- do not
%  "fix" them.
%
%  WHAT IS STILL UNMEASURED.  The collation sample is 20% of the body; a rate
%  around 0.1 errors per page cannot be excluded by a sample this size, and
%  the remaining 56 pages have been checked by the two diffs and by their own
%  transcribing agent, but not by a second reader.  Of the 22 footnotes, those
%  on the sampled pages were read; the rest rest on their batch agent's
%  reading.  spacing.py is unreliable inside footnote type -- it missed a
%  four-word letterspaced run in the note to p. 21 -- so footnote emphasis
%  throughout was decided by eye and is the weakest class in the edition.
% ============================================================
\documentclass[12pt,a4paper]{book}
\usepackage[utf8]{inputenc}
\usepackage[T1]{fontenc}
\usepackage{libertinus}
\usepackage{libertinust1math}
\usepackage{graphicx}    % for \reflectbox (the old Danish »ɔ:« id-est mark)
\DeclareUnicodeCharacter{0254}{\reflectbox{c}}  % ɔ (U+0254) = reversed c
\usepackage[danish]{babel}
\usepackage[protrusion=true,expansion=true]{microtype}
\usepackage{setspace}
\usepackage[margin=1.2in]{geometry}
\usepackage{enumitem}
\usepackage{fancyhdr}
\setlength{\headheight}{28pt}
\addtolength{\topmargin}{-13.5pt}

% Footnotes as the 1866 printing sets them: *) and **), not 1, 2, 3.
% This book is NOT like Evangelietroen og Theologien, where every note was *)
% and a constant \thefootnote sufficed: here several printed pages carry two
% notes, and the second is **).  Neither \fnsymbol (math mode: sets U+2217
% ASTERISK OPERATOR, not a text asterisk) nor footmisc[perpage] (resets on the
% TYPESET page, not the printed one) is right.  Instead the mark is set
% explicitly at each site with \fnmark{*)} or \fnmark{**)}, which restores the
% default afterwards so a missed one cannot silently propagate.
\renewcommand{\thefootnote}{*)}
\newcommand{\fnmark}[2]{\renewcommand{\thefootnote}{#1}\footnote{#2}%
  \renewcommand{\thefootnote}{*)}}

\usepackage{hyperref}
\hypersetup{
  pdftitle={Dualismen i vor nyeste Philosophie},
  pdfauthor={Georg Brandes},
  hidelinks
}

\pagestyle{fancy}
\fancyhf{}
\fancyhead[LE,RO]{\thepage}
\fancyhead[RE]{\textit{Dualismen i vor nyeste Philosophie}}
\fancyhead[LO]{\textit{\leftmark}}

\setstretch{1.3}

\title{\textbf{Dualismen}\\[6pt]
  \large i vor nyeste Philosophie}
\author{G.\ Brandes}
\date{Kjøbenhavn: Forlagt af den Gyldendalske Boghandel (F.\ Hegel), 1866}

% ---- original-page markers (paginate.py) ----
% \opage{N} is written BY HAND, at the word printed page N begins on, by
% someone who read the page.  \apage{N} is reserved for markers placed
% afterwards by machine alignment; never write it by hand.
\usepackage{marginnote}
\newcommand{\opage}[1]{\marginnote{\footnotesize\textit{[#1]}}}
\newcommand{\apage}[1]{\marginnote{\footnotesize\textit{[#1]}}}
\begin{document}
\maketitle
\thispagestyle{empty}
\newpage

\tableofcontents
\newpage

%% ============================================================
%%  PRELIMINARIES (printed pp. 1--6 = PDF 5--10)
%%  Transcribed from the image at 400--900 dpi, not from either text layer:
%%  the playbook is explicit that contents-page OCR is unreliable, and the
%%  Google layer does in fact render the p. 47 head with a spurious space
%%  before the stop.
%% ============================================================
\chapter*{Titelblad og Indhold}
\addcontentsline{toc}{chapter}{Titelblad og Indhold}
\label{ch:prelim}

% --- p. 1 ---
% Half-title.  Both lines letterspaced, `Dualismen' in a larger size; a short
% centred rule below.  The ink shelfmark across the top right corner is a
% British Museum accession mark, not printed matter.
\begin{center}
\opage{1}\emph{\Large Dualismen}\\[10pt]
\emph{i vor nyeste Philosophie.}\\[14pt]
\rule{2cm}{0.4pt}
\end{center}

% --- p. 2 ---
% Blank; British Museum oval stamp only.
\opage{2}

% --- p. 3 ---
% Title page.  `Dualismen' is a large display face and is NOT letterspaced
% here, unlike the half-title.  The slanted stroke under `G. Brandes.' and the
% pen-flourish beside it are manuscript ink in this copy, not a printed rule.
\begin{center}
\opage{3}{\Large\textbf{Dualismen}}\\[10pt]
\emph{i vor nyeste Philosophie.}\\[14pt]
\textbf{Af}\\[10pt]
\emph{G.\ Brandes.}\\[16pt]
% Publisher's device: a buckled garter enclosing an intertwined script
% monogram, the garter lettered clockwise `1770 · SG · JD · FH ·' — the
% Gyldendal founding year and the initials of its successive proprietors.
% The inner monogram is two intertwined script capitals: read as GB, but the
% left letter could equally be a looped S (i.e. SG, Søren Gyldendal), and at
% 600 dpi bitonal the two cannot be separated.  DOUBTFUL — not resolved.
[Forlagsmærke]\\[16pt]
\emph{\textbf{Kjøbenhavn.}}\\[4pt]
Forlagt af den Gyldendalske Boghandel (F.\ Hegel).\\[4pt]
{\small\emph{I.\ Cohens Bogtrykkeri.}}\\[6pt]
1866.
\end{center}

% --- p. 4 ---
% Blank; British Museum oval stamp only.  No printer's imprint on the verso,
% no copyright notice, no dedication.
\opage{4}

% --- p. 5 ---
% Indhold.  The entries carry NO chapter numbers, though the chapter heads in
% the body do, and they end in NO full stop, though the heads do.  Both
% disagreements are recorded rather than normalised; see the chapter heads
% below.  Leaders are spaced full stops and their number carries no meaning.
\begin{center}
\opage{5}\emph{\textbf{Indhold.}}\\[4pt]
\rule{1.5cm}{0.4pt}
\end{center}

\begin{flushleft}
\hfill\textbf{Side}\\[2pt]
Indledning \dotfill 7\\
Den absolute Uensartethed \dotfill 25\\
Det Antirationelle \dotfill 47\\
Dualismens Consequenser \dotfill 55
\end{flushleft}

\begin{center}
\rule{1.5cm}{0.4pt}
\end{center}

% --- p. 6 ---
% Blank.  No stamp, no printed matter.
\opage{6}

%% ============================================================
%%  1. INDLEDNING  (printed pp. 7--24 = PDF 11--28)
%%  Head as printed: `1.' on its own line, then `Indledning.' — both bold and
%%  centred, the title NOT letterspaced; the page drops its folio; the first
%%  word opens with an enlarged capital.
%%  Indhold reads `Indledning', without number and without stop.
%% ============================================================
\chapter*{1.\ Indledning.}
\addcontentsline{toc}{chapter}{1. Indledning}
\markboth{Indledning}{Indledning}
\label{ch:1}

% --- p. 7 ---
% Chapter opening: the head `1.' / `Indledning.' is already set in the
% skeleton and is NOT repeated here.  p. 7 drops its folio (first printed
% folio in the whole book is 8, on the next leaf).  The first word opens with
% an enlarged capital D — a printer's device, not emphasis, so no markup.
% No letterspacing anywhere on pp. 7--15: spacing.py flagged `Der' (p. 7,
% 1.21), `have' (p. 8, 1.18) and `fri' (p. 11, 2.45), and all three were read
% at 600--900 dpi and are ordinarily set.  No footnote on pp. 7, 8, 9, 13, 14.
\opage{7}Der er i Prof.\ R.\ Nielsens Opfattelse af Forholdet mellem
Videnskaben og Bibeltroen Noget, der nødvendigt maa vinde ham Manges
Hjerter. Man fornemmer en med Tilfredshed forenet Beroligelse, en glædelig
Bekræftelse paa, hvad man maaskee uden meget Haab har ønsket, ved at høre en
Tænker udtale sig saa strengt om den menneskelige Tankes Afmagt, ved at see
en saa stor og omfattende philosophisk Begavelse være sig Philosophiens
Begrændsning saa bevidst. Man kom til Philosophen kun altfor vant til den
indtil det Trivielle gjentagne Phrase, at den Dannede --- undertiden med
Stolthed, undertiden med Smerte --- føler sig løsreven fra de
Forudsætninger, hvorpaa den Ulærde og Enfoldige med Fortrøstning bygger, og
man gaaer bort, bestyrket i sin Tillid til, at det kun er den overfladiske og
umodne Philosophie, der fører bort fra Bibeltroen, men ikke den sande, dybe,
virkelige Forstaaelse af, hvad der kan forstaaes. Og i Tilfredsstillelsen
blander sig Beundring. Det er den samme Mand, der med saa megen Iver og Varme
har hævdet Fagviden%
% --- p. 8 ---
% First page in the book to carry a printed folio.
\opage{8}skabernes Selvstændighed imod en overmodig Speculations
Anmasselser, som staaende paa Høidepunktet af sin Tids Videnskabelighed,
iført sin Genialitets og sin Dannelses Styrkebelte har Maadehold,
Besindighed, Ydmyghed og Pietet nok til at erklære, at Livet er Mere end
Betragtningen, Charakteren Mere end Intelligentsen, Praxis Mere end Theorie
og Religionen høiere end Videnskaben. Andre Tænkere, ringere Tænkere,
fremmede, især tydske, Philosopher have ladet sig forlede til det ilde
berygtede moderne Fritænkeri, have vakt Forargelse ved Angreb paa det, som
Aarhundreders Tradition har stemplet som helligt: vor egen Philosophie vil
% The colon after `helligt' is on the page at 600 dpi and in the KB layer;
% the Google layer drops it.
under denne Tænkers Auspicier gaae i et bedre Spor; uden at give Afkald paa
den Frihed, der er Tankens Ret og Livsbetingelse, vil den, ren for de
Udskeielser, hvortil Videnskaben ellers overalt i Europa og tildels ogsaa her
i Danmark har hengivet sig, bevare Ærefrygten for det Hellige i værdig
Selvbeherskelse. Hvo har ikke nok af den golde Polemik, af
Nedbrydelseslysten og Tilintetgjørelsesraseriet, hvad er smukkere og bedre
end denne saa oprigtigt og kraftigt ytrede Bestræbelse for dog endelig engang
at stifte Fred mellem Magter, der hver for sig af uendelig Betydning
fredeligt burde samvirke til et fælles Maal!

I disse eller lignende Udtryk vil Beundringen klæde sig, og den Kreds, der
taler saaledes, er neppe lille. Thi denne Philosophie er i mere end een
Henseende en Forsoningens Philosophie. Den forsoner
% --- p. 9 ---
\opage{9}ikke blot Principer, men Partier, den søger sine Venner i modsatte
Leire: Har den end, hvad det Religiøse angaaer, sit hele Idee-Indhold fra
Kierkegaard, saa er den dog hævet over, hvad den maa opfatte som en
Indskrænkning bos dette store Genie, den seer
% PRINTER'S ERROR, transcribed as printed: `bos' for `hos'.  Verified at
% 900 dpi in the Google (British Museum) copy and at 400 dpi in the KB copy —
% both print a closed bowl, i.e. b, not the open shoulder of h.  Both text
% layers independently read `bos'.  The error is in the forme, not the copy.
ikke alene som næsten Alle det i mange Maader Vækkende og Gode i det Liv, der
er udgaaet fra Grundtvig og hans Tilhængere, men den har Øie for den dybe
Sandhed i Grundtvigs Opdagelse, som Kierkegaard og saa Mange med ham ere ude
af Stand til at finde Synspunktet for. Thi denne Philosophie forstaaer at
vurdere den friske Umiddelbarhed saa godt som den gjennemførte Reflexion; kun
Mellemsphæren, kun den uklare Sammenblanding af det Videnskabelige og det
Religiøse, med eet Ord Theologien er den fjendsk; og dog --- den er forvisset
om, at Theologerne, naar de ret faae Syn paa, hvad i høieste Forstand er
deres egen Interesse, ville føle Trangen til at samles under dens brede
Vinger med deres forhenværende Modstandere; thi hvad de bedste iblandt disse
ville, det vil ogsaa den: det ubeskaarne uantastelige Bibelords Hellighed,
Troens Uafhængighed af den stadigt mere fordringsfulde og hensynsløse moderne
Videnskab. Den indeslutter sig endelig ikke med sine Indviede i et afspærret
Rum, den er betænkt paa at skaffe sig Udbredelse til de almindeligt Dannede,
ja til Folkets store lærvillige Masse, som det er en Lyst at undervise for
den, der er lykkelig nok til at have Evnen dertil, den er ligesaa populær,
% --- p. 10 ---
\opage{10}naar den vender sig udad, som den er streng og exact i sit Forhold
indadtil, ligesaa blød og forsonlig, naar det gjælder om at forene, hvad
Stort og Udmærket vort Fædreland har frembragt i gudelig og kirkelig
Retning\fnmark{*)}{Se f.\ Ex.\ Nord.\ Universitetstidskr.\ Tillægshefte til
10de Aarg.\ 1866 p.\ 75.}, som den er charakterfast og ubønhørlig imod
% Printed mark: *) — one note only on this page.  `Se', single e, against
% `See' in the note on p. 15; both as printed.
Theologiens (for Almenheden ligegyldige) Formidlingsforsøg.

Hvor meget er der ikke i et saadant Program, der tiltaler det Bedste i
Mennesket, og som til alle Tider vilde gjøre Indtryk paa enhver moden og
udviklet Aand, og hvor Meget er her endelig ikke, der maa ramme netop vor
Tids og vort Lands Mennesker!

Det 19de Aarhundrede gjorde tidlig Front mod det 18des Philosophie; men
hertillands maaskee mere afgjort, i hvert Fald mere udholdende, end
nogensteds, og den danske Litteratur bærer, hvor heftigt vore Theologer end
have befeidet hinanden, hverken mange eller stærke Spor af den kritiske
Bevægelse, der i indeværende Aarhundrede charakteriserer den tydske og
franske Videnskab. Tiden for de religiøse Bevægelsers Gjenfødelse har Intet
tilovers for Rationalismen, den afskyer den gamle Rationalismes Fladhed, den
tvivler høiligt paa, at der kan komme noget Godt ud af paany at rokke ved det
positivt Religiøses Autoritet, den seer nødig Opgivelsen af nogetsomhelst af
Aabenbaringens Udenværker; thi den føler Nødvendigheden af „principiis
obstare.“ Denne
% --- p. 11 ---
\opage{11}Aandsretning kommer Prof.\ Nielsen imøde. Han sikrer det Positives
alleryderste Forskandsning, idet han anbringer et gabende Svælg imellem
Videnskaben og Religionen, og kaster Theologerne paa Hovedet ned deri.

Men det 19de Aarhundrede er tillige Naturvidenskabens og de store
Opdagelsers Tidsalder. Er den klar over Noget, saa er det over det, at paa
Videnskabens Omraade maae vi have Marken fri, maae vi kunne forske og
bestemme uanfægtet af en forældet Orthodoxies forudfattede og indskrænkede
Domme. Denne Bevidsthed kommer Prof.\ Nielsen imøde. Naar Talen er om
Videnskab, da er Talen ikke om Tro, da staaer Prof.\ --- eller vil
idetmindste staae --- fri og uhindret af enhversomhelst bindende Fordom,
ingen nok saa utaalmodig Fritænker har herom kunnet udtale sig heftigere end
han.

Hvad vil man da Mere, hvad kan man endnu ønske udover denne storartede
Gjennemførelse af suum cuique! hvo vil vel her gjøre Indsigelse, med
Undtagelse maaskee af Efternøleren, der endnu ikke forstaaer, at Dogmatikens
Tid er forbi, eller af den Bornerte, der ikke kan taale, at den store
Almenhed delagtiggjøres i Philosophiens Resultater, eller den Enkelte, der
optræder ikke for Sagens Skyld, men for at bringe sit eget lille Jeg i
Forgrunden ved Hjælp af den store Sag\fnmark{*)}{Jfr.\ R.\ Schmidt: Om
Selvmodsigelsen i Prof.\ Nielsens Lære pag.\ 44.\endgraf
Ligesom Hr.\ Schmidt antager, at der maa ligge noget moralsk Graverende til
Grund for den Handling at træde op imod Prof.\ Nielsen, saaledes kan Hr.\
Henrik Scharling ikke fatte, at man udtaler sig mod Theologien af andre
Grunde end personlige. Han erklærer, at jeg synes at staae i den Formening,
at naar Theologien ikke kan gjøre mig til en Christen, saa er den ufrugtbar
og upraktisk og tilføier belærende, at man naturligvis ligesaa lidt kan blive
en Christen alene ved at studere Theologie, som man kan gjenvinde sin Helbred
alene ved at studere Medicin. Skjøndt Hr.\ Scharlings Morsomhed som
Litteraturens Mentor ikke er ubekjendt, og skjøndt man jo vel maa være
forberedt paa at træffe nogen Familiaritet med Alvidenheden i et theologisk
Ugeskrift, var dette Udbrud af Hr.\ Scharling mig dog en kjær Overraskelse.
Det er saa behageligt i et Tidsskrift at finde Oplysninger til sin egen
Sjælehistorie og tilmed Oplysninger, som man tilforn var ganske uvidende
om.}?
% Printed mark: *).  THIS NOTE RUNS OVER ONTO PRINTED p. 12: the first
% paragraph (`Jfr. R. Schmidt … pag. 44.') and the opening of the second
% (`Ligesom Hr. Schmidt antager, at der maa ligge') stand at the foot of
% p. 11; the second paragraph continues, unmarked and below a fresh rule, at
% the foot of p. 12 from `noget moralsk Graverende'.  Set here as ONE note at
% its mark on p. 11, per BATCH-AGENT.md.  `\endgraf' rather than `\par'
% because \fnmark is not \long and a literal \par token in its argument
% aborts the compile.  `naturligvis' (not `naturligviis') is as printed —
% confirmed at the image, and both layers agree.

% --- p. 12 ---
\opage{12}Lad os engang vende Medaillen og see den paa den anden Side. Denne
Lære, der tiltaler saa meget Godt i Menneskene og viser sig saa
tidssvarende, skulde den mon ikke ogsaa henvende sig til noget mindre Godt,
til det Lavere, ja til det Laveste i os? Vi Mennesker ere af Naturen
tilbøielige til Slaphed, Delthed, Adsplittethed. De Fleste have gjerne een
Livsanskuelse i Glæde, en anden i Sorg, een i Hjemmet, en anden i Kirken.
Læreren hylder helst eet Princip, naar han taler, et andet, naar han handler,
Præsten har Møie med at leve som han prædiker, Kjøbmanden har gjerne een
Moral i Forretninger, en anden i Livet, og Digteren besynger tidt og ofte
Ide%
% --- p. 13 ---
\opage{13}aler, som allermindst gjenfindes i hans egne Handlinger. Den ene
Sphære kommer jo ikke den anden ved, den ene Haand behøver jo ikke at vide,
hvad den anden gjør. Og nu Videnskabsmanden! Opdragen i en vis bestemt
kirkelig Tro, fra Barndommen af paavirket af en vis bestemt tilfældig
Forestillingskreds, den, som hersker i hans Kirke, i hans Land og i den By,
hvori han lever, møder han i Videnskaben Facta, der staae i Strid med denne
Tro. Her er en Kamp at udstride, en Seir at vinde. Men kunde Freden ikke
faaes for bedre Kjøb? Forfærdelige Tanke, at den Tro der prædikes i Ebeltoft
eller i Paris eller i St.\ Petersborg ikke skulde være den høieste Sandhed!
Forfærdelige Utroskab, at bryde med en Tro, til hvilken Ens ebeltoftske eller
parisiske eller St.\ Petersborgske Forældre satte al deres Lid! Hvad er da
naturligere end den praktiske Udvei at være En i Videnskaben, en Anden i sit
Liv, at leve Livet igjennem i en chaotisk Verdensanskuelse med Brokker af
Fortidens og Stumper af Nutidens Livsbetragtning, men dog i sine Skrifter at
fremtræde iført det 19de Aarhundredes videnskabelige Drapperi.

Dog saa naturligt dette end falder, der er noget Pinligt i denne Halvhed.
Man samler sig vel nødigt til den kraftige, enkelte Afgjørelse, men man lider
heller ikke at vakle principløst mellem modsatte Synsmaader. Ingen Accord!
det er jo Løsenet i vore Dage, saa langt heller en dygtig Disharmoni!
Halveres vil man, naar det ikke kan være anderledes, men saa
% --- p. 14 ---
\opage{14}vil man ogsaa halveres tilgavns. Hvilken Lindring for den
bespændte Sjæl giver nu ikke den nye Lære, der vel nødtvungen indrømmer
Videnskaben dens Ret, men ligefuldt lader mig min Barnetro, ja om det saa
skulde være min Ammestuelærdom i god Behold, som kort sagt hæver
Principløsheden selv til Princip og skjænker den Usikre netop hvad han
attraaer, at maatte blive hel i Halvheden, at maatte aande fuldt ud som et
dobbelt Menneske. Hans Forstand trænger til Luft: den faaer et Spillerum, en
Tumleplads uden Grændser, hans Følelse, der ligger i Splid med Forstanden,
hans Villie, der i Trang til udvortes Støtte, vanskeligt slipper nogen fra
Fædrene overlevet Troessætning, behøver intet Afkald at gjøre. Troen er
berettiget, saasnart den blot er fast, dens Rige er uden Skranker, dens
Omraade har ingen Grændse tilfælles med Erkjendelsens Land. Hvert af de
tvende Riger har sine Love, sin Regent, ja sin Gud, de kunne ikke splidagtigt
strides om det Sande; thi selve denne dybe Splid imellem dem, det er
% NOT a defect of the printing: the Google (British Museum) copy shows a
% loose space inside `selve' (`thi s elve'), but the KB copy at the same line
% is normally spaced, so it is a slipped space in that one impression.  Both
% text layers read `selve'.  Transcribed `selve'.
Sandheden, det er Freden. Som naar tvende Stormagter staae fuldtrustede imod
hinanden, men Mægling indtræder, ingen Krig bryder ud, og der stiftes et
Forlig, i hvilket hver af Parterne fuldtbevæbnet forbeholder sig alle sine
Rettigheder, saaledes kvæles her Debatten under en uveirssvanger Pause: ingen
af de to Modstandere opgiver en Tøddel. Betragtet som en midlertidig
Situation eller som en Vaabenhvile efter fordums Kampe kan denne Stilstand
have sin Nytte, den er endog af Værdi for Viden%
% --- p. 15 ---
\opage{15}skaben, hvis denne virkelig faaer Lov til at leve sit Liv
uforstyrret af fremmed Indblanding; men skal den opfattes som en Afgjørelse
af Striden, som den længe ventede Løsning af et stort Problem --- og for
dette bliver den udgivet\fnmark{*)}{See f.\ Ex.\ „Fædrelandet“ 1866 Nr.\
172.} --- da er dens Svaghed saa iøinefaldende, at det vilde være
% Printed mark: *).  `See', double e, as printed — cf.\ `Se' on p. 10.
beskjæmmende for vor Videnskab, om den ikke blev paapeget.

Vi saa, hvad der kunde friste endogsaa Videnskabsmanden til at slutte sig til
den nye Lære og for sin Part realisere den. Men kan Fristelsen blive selv ham
for stærk, hvor meget større maa den da ikke være for Lægmanden. Sæt han fra
sin Barndom af har en Trang til ubetinget Hengivelse og inderlig Tro. Sæt
saa, at, naar hans Hjerte søger den positive Religion, hans Fornuft samtidig
frastødes af mangen med Religionen, som det synes, uopløselig forbunden
Cruditet. Forkaster han den nu ubetinget, da forsmægter hans Hjerte bestandig
under sit Savn, antager han den derimod uden Forbehold i dens overleverede
Skikkelse, da er det, som sloges der ham et Jernbaand om Panden. Saaledes
sidder han fast i Valgets Dilemma. Sæt saa, at han en Aften træder ind i en
Høresal, hvor Tilhørere af begge Kjøn, Dannede og Udannede, have trængt sig
sammen. Det Foredrag, han kommer til at høre, er rettet mod den speculative
Theologie. Med al den Sikkerhed, som Lærdom, logisk Virtuositet og livfuld
Veltalenhed giver, fremstiller Taleren,
% p. 15 ends mid-sentence; `fremstiller Taleren,' continues on p. 16, which
% belongs to the next batch.
% --- p. 16 ---
% Continues the sentence broken at the foot of p. 15 (`fremstiller Taleren,').
% No footnote on p. 16.  No letterspacing: spacing.py flagged `man' (1.18,
% `da man har lært'), read at 700 dpi and ordinarily set.
% The em-dashes round `naturligvis bestandig fra Videnskabens Standpunkt' are
% on the image; the Google layer renders the first as nothing and the second
% as a hyphen glued to `Usandsynligheden' (ORTHOGRAPHY §3d).
% `naturligvis', single i, as printed — p. 16 is in the eleven (ORTHOGRAPHY §1).
% Google invents a full stop after `videre, han'; KB and the image have none.
\opage{16}hvor farligt det i Troessager er at række Videnskaben, om det saa
kun var den lille Finger. Paa Videnskabens Vegne nedlægger han Protest mod de
bibelske Underes Troværdighed. Han fremstiller maaskee først --- naturligvis
bestandig fra Videnskabens Standpunkt --- Usandsynligheden, Absurditeten af
nogle af de Mirakler, der for hans Tilhøreres Phantasie kun have liden Nimbus,
men som ere deres Forstand til des mere Anstød. Endnu føler den lyttende Kreds
sig ikke foruroliget, snarere fornemmer den en Sindets Lettelse ved engang at
give den tilbagetrængte Forargelse frit Løb. Men Taleren gaaer videre, han
nærmer sig maaskee Mirakler, der for Alle staae som hellige Mysterier, det
synes et Øieblik, som vilde han forgribe sig paa dem, og en Bæven gjennemfarer
Forsamlingen. Consequensen af Theologernes Forklaringsforsøg er nu gjort
tydelig. Da vender Taleren Foredragets Timeglas. Med varm Begeistring, med
endnu større Veltalenhed fremstiller han Miraklet seet i Troens Lys. Og nu, da
den kritiske Tilbøielighed har faaet Lov at udrase, nu, da man har lært at
skjælve for den Tankegang, man nylig var villig til at slaae ind paa, nu
omfatter man med samme Lidenskab Stort og Smaat, det mest ophøiede religiøse
Bud og det mindst betydende Under, Alt synes lige helligt og kostbart, da
Faren for at miste det neppe er forvunden. Hvad har nu Tilhøreren i en saadan
Time ikke gjennemfølt og gjennemlevet, hvor have alle hans Evner tumlet sig
frit! Fornuften og Troestrangen ere lige tilfreds\-%
% --- p. 17 ---
% Printer's signature `2' stands alone at the foot of this page, below the
% footnote: it is not text and is not transcribed.
% One footnote, mark *).  No letterspacing on the page or in the note —
% checked at 700 dpi, the note included, because spacing.py's fit is unreliable
% in the smaller footnote type (see p. 21).
% `Propædeutik' with the ligature, as printed and as Google reads it here;
% ORTHOGRAPHY §3b records Google losing the same æ in the notes to pp. 55, 59.
\opage{17}stillede, hin fordi den har faaet Tilladelse til at kaste hele
Mirakelflokken ud af Videnskabens Forport, denne, fordi den har faaet
fuldkommen Frihed til at lukke hele den samme Skare ind ad den Bagport, som
aabnes den i Individets praktiske Existens. Frem og tilbage synes lige langt;
men den Hørende har brugt sine aandelige Kræfter, og han har een Gang for alle
lært, hvorledes han har at behandle Forstanden, hver Gang dens Utaalmodighed
paany skulde udsætte ham for Forargelsens Anfægtelser. I hans Glæde undgaaer
det ham, at hvad her er foretaget kun er en fingeret Bevægelse, at, da det
Gudsbegreb og det Begreb om Verden, som Prof.\ N.\ stiller imod Theologiens,
efter Prof.'s egen Opfattelse ikke er det sande, ikke er det, som har Realitet,
saa er den Strid, som her er ført mod den speculative Theologie kun en
Skinkamp, en kvantvis Fægtning, som Holberg kalder det, og Alt, hvad her er
skeet, kun en Ompostering af det Suprarationale fra Videnskabens til Livets
Omraade, en taktisk Forandring, analog med den, som vilde være indtraadt i
Chemien, hvis Chemikerne i sin Tid vare blevne enige om, ikke mere at taale
Phlogiston i deres Videnskab, men kun for udenfor Videnskaben at være des mere
glødende forvissede om dets Existens.

Eller den Famlende og Usikre tyer maaskee til Litteraturen og standser en
skjøn Dag ved Læsningen af en Pagina som denne\fnmark{*)}{R.\ Nielsen:
Forelæsninger over Philosoph.\ Propædeutik. Efteraarshalvaar 1865 pag.\ 159.},
% Printed mark: *) — one note only on this page.  It sits between `denne' and
% the following `)', i.e. `denne*),'; the comma is outside the mark.
der indeholder Slutningen
% --- p. 18 ---
% PRINTER'S ERROR, transcribed as printed: the letterspaced word reads
% `teleolo-' at the line end and `logisk,' at the head of the next, i.e.
% `teleolologisk' — the syllable `lo' is set twice.  Read at 700 dpi on both
% lines.  Both text layers independently give the same doubled form, and the
% word is set correctly (`teleologisk', and NOT letterspaced) fourteen lines
% below in C.'s second speech, which is what shows the first to be a slip of
% the compositor's.  Not corrected.
% The letterspacing of `teleolologisk' is on the image and in KB's spurious
% internal spaces (`t e l e o l o - l o g i s k'); spacing.py flagged both
% halves.  It is the only emphasis on the page.
% The quotation closes with THREE spaced full stops (counted at 700 dpi),
% against the six on p. 24.
\opage{18}af en Samtale, hvori Theologiens Forsvarer kommer ynkeligt tilkort.
Theologen C.\ har opstillet den Sætning, at i objectiv Betydning er Alt eller
Intet at kalde Under. Hertil bemærker saa A.\ som følger:

„At et Menneske vandrer paa Havet uden at synke, er da ligesaa stemmende med
Naturlovenes Uforanderlighed, som at En falder i Vandet og drukner? C.\
Mennesket, der ellers er underkastet Tyngdens Love, kan meget vel vandre paa
Havet, naar det skeer \emph{teleolologisk}, ifølge en fortsat Skabelse! A.\ Men
saa har det, hvad Naturlovenes Uforanderlighed angaaer, vel heller ikke, naar
vi ret besinde os, noget for den dybere Tænkning Anstødeligt, at en Pave, Leo
XII.\ har kunnet beatificere Minoriten Julianus, som i vort vantro,
revolutionære nittende Aarhundrede fik stegte Fugle til at flyve? C.\ At jeg
for mit Vedkommende antager Underet med de stegte Fugle for en dum Fabel, en
Præsteløgn, hidrører slet ikke fra, at jeg i Begivenheden selv seer noget Brud
paa Naturlovenes Uforanderlighed; men simpelthen derfra, at jeg ikke er
Catholik og følgelig ikke kan overbevise mig om, at Begivenheden er skeet
teleologisk, ifølge en fortsat Skabelse. Var jeg Catholik og havde en Catholiks
Tro, da skulde jeg vel være istand til at forsone min Tro med min Viden, mit
Skabelsesbegreb med mit Naturbegreb . . .“

Hvad er her nu skeet? Idet Theologiens Forklaringsforsøg ere latterliggjorte,
idet den indbildte Forstaaelse af Underet har vist sig at være indbildt,
% --- p. 19 ---
% `Ælde' with the capital ligature, READ AT THE IMAGE at 700 dpi — the Google
% layer gives `Elde' (ORTHOGRAPHY §3a, the p. 19 entry); KB has `Æ lde'.
% Every other word on pp. 16--24 opening with a capital vowel was checked at
% the image for the same fault, and none is an Æ.
% `Del', single e, as printed — p. 19 is in the seven (ORTHOGRAPHY §2).
% `din' is letterspaced: confirmed at 700 dpi, flagged by spacing.py, and
% split as `d in' by KB.  The only emphasis on the page.
% Google reads `sympatketiske'/`bar sin Grund' in KB; the image has
% `sympathetiske' and `har sin Grund'.
% Printer's signature `2*' stands alone at the foot: not text.
% NB the reported speech opened here at `At jeg for mit Vedkommende' carries
% NO opening „ on the page — checked at 700 dpi — though it is closed with “
% on p. 20 at `Almagt.“'.  Transcribed as printed; see the note there.
\opage{19}rammer Spotten ikke alene den speculative Dogmatik, som den der ---
principielt --- formaaer at begribe et saa afsindigt Under, men som den, der
ikke er istand til at afvise det. Saaledes falder der da et latterligt Skjær
over selve den miraculøse Verden, som Dogmatiken behandler, den Verden nemlig,
hvori stegte Fugle kunne bringes til at flyve. Atter her føler den
sympathetiske Læser sin Forstand tilfredsstillet; skjøndt Intet ved sin Ælde
ærværdigt Mirakel er bleven bragt endog blot under Skyggen af en
Forstandskritik, saa gotter man sig dog ved Tanken om, at man for sin Del ikke
er nødsaget til at tye til saadanne Argumenter som C.\ for at værge sig mod
Antagelsen af den Art Mirakler, som det med Fuglene. Men denne Følelse har sin
Grund i ren Illusion. Thi optager man Samtalens Traad og retter Repliken mod
Hr.\ A., da vil det vel vise sig, at hans Argumenter ikke ere et Haar bedre end
hans ydmygede Modstanders, og at hans Situation er nok saa besynderlig som
Hr.\ C.'s. Spørger man ham nemlig blot med nogle Ords Forandring saaledes: Har
det da, hvad Naturlovenes Uforanderlighed angaaer, for \emph{din} inderste
Bevidsthed noget Anstødeligt, at en Mand i vort Aarhundrede fik stegte Fugle
til at flyve? da maatte Svaret lyde som saa: At jeg for mit Vedkommende antager
Underet for en dum Fabel, en Præsteløgn, hidrører slet ikke fra, at jeg i
Begivenheden seer et Brud paa Naturlovenes Uforanderlighed, thi indeholder den
end et saadant, tør jeg dog ikke be\-%
% --- p. 20 ---
% UNPAIRED CLOSING QUOTE, transcribed as printed: `Almagt.“' closes a reported
% speech that was never opened with „ on p. 19.  Read at 700 dpi on both
% leaves.  This is the first of the two places in this batch where the quote
% count drifts off zero on purpose; the second is on p. 24.  Batch total:
% 7 opening „ against 9 closing “ — the drift is exactly these two.
% No footnote, no letterspacing (spacing.py: none; confirmed at the image).
% `venter vel ikke, at det vil skee' — comma, as Google and the image have it;
% KB reads a semicolon and is wrong.  `Fortællinger' (KB: `Fortælling'),
% `Regnbuens Fremkomst' (KB: `i remkomst') — both settled at the image.
\opage{20}negte dets Mulighed, men simpelthen deraf, at jeg ikke er Catholik og
følgelig ikke troer, at der er skeet Mirakler paa Jorden siden det første
Seculum; derfor brugte jeg Ordene: „i vort vantro, revolutionære nittende
Aarhundrede“. Var jeg en Catholik, og havde en Catholiks Tro, da skulde jeg vel
være istand til at forsone min Tro med min Viden; thi jeg staaer ikke min
Fornuft til Regnskab for, hvad jeg troende antager. Naturlovenes
Uforanderlighed er en Norm, som kun staaer fast i en Naturforfatning, der
forefindes i Videnskaben og paa Papiret, men som ikke har hjemme i den
virkelige Verdensforfatning, hvortil jeg i min Existens forholder mig som en
lydig Undersaat. Mod den strider intet Mirakel, som Troen erklærer for helligt
og som styrker Forvisningen om Altets absolute Afhængighed af den ubetingede
Almagt.“ Med andre Ord: ogsaa Hr.\ A.\ har den Mening, at i og for sig er
Naturlovenes Uforanderlighed ikke til Hinder for, at de stegte Duer kunne flyve
os ind i Munden; han venter vel ikke, at det vil skee, men kun fordi den Tid,
hvori Sligt hændtes, er omme.

Naar man til En af vore Ny-Orthodoxe, der ikke ville indrømme Fornuft eller
Videnskab den svageste Ret til at tale med om Religionens Gjenstand, bemærker,
at det dog i vore Dage ikke gaaer an at fastholde Fortællinger som den om
Sprogenes Oprindelse fra Guddommens Vrede over en formastelig Taarnbygning
eller om Regnbuens Fremkomst som corporlig Ting i Anledning af en bestemt
historisk
% --- p. 21 ---
% One footnote, mark *), set after `afgjøre' and before the full stop:
% `at afgjøre*).' — read at 700 dpi.
% TWO letterspaced passages, and spacing.py found only the first:
%   body   — `vil' in `Hvad Bibelen vil lære os' (flagged; confirmed at
%            700 dpi, and split as `v il' by KB);
%   NOTE   — `Troen, Religionen, Videnskaben, Villien', a run of four nouns
%            with their commas, NOT flagged by spacing.py, which does not fit
%            the smaller footnote type reliably.  Confirmed at 700 dpi against
%            `taler om' and `som om personlige Væsener' on the same two lines,
%            which are tight.  KB's `T ro en , R e lig io n e n , V idenskaben,
%            Villien' is the corroborating signal (ORTHOGRAPHY §3g).
%   NOT letterspaced, though KB splits them: `Udtryk', `kun et blandt',
%   `Vidnesbyrd', the title `Indledning til Ethiken' (quoted only), and
%   `Dyden er ingen Fersken' — all four read at 700 dpi and ordinarily set.
%   This is ORTHOGRAPHY §3g-bis: KB's splits inside a footnote are mostly noise.
% `Natur- eller Sprogvidenskab' keeps a REAL hyphen — suspended compounding,
% not a line-break hyphen.
% `beroer' (KB: `heroer'), `skjønne' (KB: `skjønue') — settled at the image.
\opage{21}eller uhistorisk Begivenhed uden som skjønne og poetiske Sagn, saa
svarer den Paagjældende i Almindelighed, at Bibelen slet ikke vil lære os nogen
Natur- eller Sprogvidenskab, overhovedet slet ikke vil berige vor Viden, at
dette er en usalig Misforstaaelse o.\ s.\ v. Dette Svar beroer tydeligt nok
paa en Feiltagelse. Hvad her siges om Bibelen gjælder ikke den, men den
religiøse Forkyndelse, det sædelige Bud. Dette henvender sig ikke til vor
Videbegjærlighed, vil ikke bibringe os nogen Kundskab om Sprogenes eller
Naturphænomenernes Oprindelse. Hvad Bibelen \emph{vil} lære os, er vanskeligt
at afgjøre\fnmark{*)}{At et saadant Udtryk hyppigt høres, er kun et blandt
mange Vidnesbyrd om den høist ucorrecte Sprogbrug, der har indsneget sig hos
os, i Kraft af hvilken man bl.\ A.\ taler om \emph{Troen, Religionen,
Videnskaben, Villien} som om personlige Væsener. Exempler paa denne
Metaphorphilosophie yder Dr.\ Heegaards „Indledning til Ethiken“ i rigeste
Maal. Her blot eet: Naar han pag.\ 446 erklærer, at Villiens Concentration ud
over det Rationelle er et Skridt, hvorved Villien ikke gjør sig nogensomhelst
theoretisk Betænkelighed, da mindes man om det Gamle: Dyden er ingen Fersken.
Man kunde med samme Ret tale om Skridt, hvorved Hjertet ikke gjør sig mindste
Hovedbrud.}. Men spørger man simpelthen, om Bibelen ikke --- hvad enten den nu
vil det eller ei --- lærer os noget Saadant, da bliver Svaret Ja, og den
Orthodoxe har følgelig, hvis han opgiver Bibelens Autoritet med Hensyn til
disse Punkter, en Gang for alle brudt med sit Princip.

Det kunde nu synes, som om en Discipel af
% --- p. 22 ---
% No footnote, no letterspacing (spacing.py: none; confirmed at the image).
% `den ethiske Fordring, der sikrer' — the comma is on the image but sits low,
% dropped almost to the descender line; KB misses it, Google has it.
% Google's layer puts a lone `.' between `Ukrænke-' and `lighed': that is a
% speck of dirt in the gutter, and so is the mark between `selv' and `det'.
% `forhaabentlig' (KB splits it `forhaa bentlig' in Google, `forhaabentlig' in
% KB) and `altfor' (KB: `altfoi') — settled at the image.
% `Prof. N.s Lære' — no apostrophe, as printed, against `Prof.'s' on p. 17.
\opage{22}Prof.\ Nielsen i dette Tilfælde ganske vilde kunne undgaae de
Vanskeligheder, hvori den Orthodoxe her finder sig hildet. Thi Ingen har jo
stærkere end Theologiens utrættelige Bekjæmper gjort Adskillelse imellem det,
der vedkommer den religiøse Bevidsthed og det, der ikke vedkommer den. Man
skulde troe, at disse Facta som betydningsløse for det religiøse Liv efter
Prof.\ N.s Lære maatte hjemfalde til Videnskabens Kritik og saaledes tabe deres
exceptionelle Stilling. Men nei, det viser sig, at de opfattes som hellige
Kjendsgjerninger. De bevares i fuld Ukrænkelighed som Mumier af den selv det
Døde ikke opgivende Tro, og det er, saa uforstaaeligt det end lyder, det Ethiske
i Mennesket, den ethiske Fordring, der sikrer disse Lig mod den kritiske
Opløsning: den ethiske Grundfordring nødvendiggjør et ethisk Ideal; det Gode
kan ikke existere uden den Gode, den absolut Gode atter kun som almægtig, den
Almægtige kun som absolut skabende, Skabelsen kan atter kun være foregaaet i de
sex Dage, (v.\ Heegaard anf.\ St.\ p.\ 392) og hermed ere de første Led knyttede
af den Kjæde, der vel maa tænkes at skulle sammenlænke alle det gamle og nye
Testamentes hellige Facta. Den er ikke udført endnu, men Paavisningen af den
hele Nexus vil forhaabentlig ikke lade altfor længe vente paa sig, og den
fuldstændige Sammenstilling vil da ligge for som et gjennemført System, hvilket
det ikke vil være raadeligt for Nogen at understaae sig til at kalde
theologisk. Spørger man altsaa Prof.
% --- p. 23 ---
% No footnote, no letterspacing (spacing.py: none; confirmed at the image).
% Four quotation marks on the page and they balance: „Et ligefrem … reale.“
% and „Men“ … „indseer Du … det ei?“.
% `nødvendigvis', single i, as printed.  `Viden til Dommer' (KB: `Yiden'),
% `uhistoriske', `Konster' — settled at the image.
\opage{23}Nielsens Discipel: Kan man forkaste det gamle Testamentes Fortælling
om den babyloniske Sprogforvirring eller det nye Testamentes om Dæmonernes
Uddrivelse og endda forholde sig til det sædelige og religiøse Ideal, da vil han
maaskee i det første Øiebliks Forbløffelse svare: Ja; men saasnart han har
sundet sig, vil han consequent benegte det: det Ethiske kan ikke undvære det
Religiøse, det Religiøse ikke Aabenbaringen, Aabenbaringen ikke disse
Forudsætninger. Sæt saa man indvendte: Forsøg dog engang at see frit og lige
paa Tingen, døm uden alle Bihensyn, antager Du disse Facta eller ei? da vilde
Svaret nødvendigvis lyde omtrent som følger: „Et ligefrem Ja eller Nei kan jeg
ikke give paa dit Spørgsmaal; thi ellers vilde jeg gjøre Viden til Dommer om
Ting, der ligge udenfor dens Competence. Som sanddru Menneske antager jeg i
Sandhedens Interesse, at disse saakaldte Facta ere ganske uhistoriske, men som
moralsk Individualitet antager jeg i det Godes Navn disse samme Facta for
fuldstændigt paalidelige og reale.“

„Men“ udbryder Spørgeren „indseer Du da ikke, at der gives en tredie simplere
Betragtningsmaade af saadanne Beretninger. Jeg spørger Dig ikke, om de ere
sande, thi med Ordet Sandhed kan der gjøres mange Konster, men svar blot paa,
om Du antager, at det, de meddele, nogensinde er skeet. Troer Du, at Regnbuen
blev sat paa Himlen, eller troer Du det ikke, troer Du at hine Afsindige vare
besatte af Dæmoner, eller troer Du det ei?“ Da lyder Svaret atter:
% --- p. 24 ---
% Last page of chapter 1.  The chapter ends with the verse; below it, a third
% of the way down an otherwise blank page, stands a short centred DOUBLE rule
% (two thin rules, about 1.7 cm wide, centred on the measure).  Nothing else
% is printed on the leaf.  p. 25 opens chapter 2 and is another batch.
% No footnote, no letterspacing anywhere on the page — the verse included,
% read at 700 dpi: it is set in the ordinary roman, in the body size, and its
% only distinguishing marks are the indentation and the surrounding space.
% SECOND UNPAIRED CLOSING QUOTE, transcribed as printed: the reported speech
% `I det Sandes Interesse …' opens with NO „ (read at 700 dpi) but closes with
% “ after the points.  Cf.\ p. 20.  Batch quote count: 7 „ against 9 “.
% SIX spaced full stops after `derimod', counted at 700 dpi by column
% projection — not three, as on p. 18.  The raised speck above the first is
% dirt on the sheet, not a seventh point.
% THE VERSE: two lines, centred neither on the measure nor on each other but
% set as a block indented about 2.9 em from the text left margin, both lines
% flush at that indent (the opening „ aligns with the `J' below it, not hung
% outside).  No attribution, no source line, no rule above or below it, and
% the second line is NOT further indented.  Set here as a `verse'.
\opage{24}I det Sandes Interesse antager jeg o.\ s.\ v., i det Godes Navn
fastholder jeg derimod . . . . . .“ og dermed maa vel Samtalen være endt. Men
her kommer selv det forsonligste Sind til at mindes det lille, skarpe Vers:

\begin{verse}
„For Aandens Sol eller Præstepraasen?\\
Ja eller Nei! kun ingen Vaasen!“
\end{verse}

\begin{center}
% The chapter-end ornament: a short centred double rule on the lower part of
% p. 24.  Not a page rule and not a footnote rule.
\rule{1.7cm}{0.4pt}\\[1.5pt]
\rule{1.7cm}{0.4pt}
\end{center}

%% ============================================================
%%  2. DEN ABSOLUTE UENSARTETHED  (printed pp. 25--46 = PDF 29--50)
%%  Head as printed: `2.' / `Den absolute Uensartethed.'
%%  Indhold reads `Den absolute Uensartethed', without number and without stop.
%% ============================================================
\chapter*{2.\ Den absolute Uensartethed.}
\addcontentsline{toc}{chapter}{2. Den absolute Uensartethed}
\markboth{Den absolute Uensartethed}{Den absolute Uensartethed}
\label{ch:2}

% --- p. 25 ---
% Chapter opening: the head `2.' / `Den absolute Uensartethed.' is already set
% in the skeleton and is NOT repeated here.  p. 25 drops its folio (first folio
% in this chapter is 26).  The first word opens with an enlarged capital D — a
% printer's device, not emphasis, so no markup.  The head is bold and is NOT
% letterspaced: spacing.py flagged `absolute' (1.24) on this page, and that
% flag is the bold head, read at 600 dpi.  No letterspacing anywhere in the
% body of p. 25; no footnote on pp. 25--29, 31, 33, 34.
\opage{25}Det træffer sig saa heldigt eller saa uheldigt, alt efter det
Standpunkt, hvorpaa man befinder sig, at Prof.\ R.\ Nielsen i Spidsen for sin
Lære om Forholdet mellem det Religiøse og det Videnskabelige opstiller en
Grundsætning, hvorpaa den hele Bygning hviler. Man behøver da ikke
møisommeligt at opspore Theoriens Grundtanke, den er udtrykt saa bestemt, at
den paa ingen Maade kan forfeiles; dersom man nu ikke kan gaae ind paa den,
saa behøver man ikke at indlade sig med det Øvrige; thi enten er hin Tanke
følgerigtigt gjennemført, og da vil den gjennemtrænge den hele Udvikling og
paa alle Punkter umuliggjøre En Tilegnelsen, eller den er ikke udfoldet ved
% `En Tilegnelsen' read at 600 dpi and confirmed in both layers.  Not an error:
% `En' is the old capitalised indefinite pronoun in the dative ("for one"), as
% in `det gjør En Livet surt'.  Left exactly as printed.
Consequens, og saa kan man jo endnu langt mindre ønske et nærmere
Bekjendtskab med alle Enkeltheder. Denne Sætning, der i adskillige Artikler i
„Fædrelandet“ er bleven betegnet som en tydelig Selvmodsigelse, lyder
saaledes: Samme Bevidsthed kan uden Modsigelse forene Tro og Viden som
absolut uensartede Principer. „Uden Modsigelse“ det er den Sikkerhedsventil,
der er anbragt for at forhindre Sætningen fra at springe
% --- p. 26 ---
% THE DENSEST EMPHASIS PAGE IN THE BATCH.  Brandes letterspaces the two general
% propositions he is about to attack, and leaves his own two interpolations
% (the ɔ: gloss, and the `skulle vi springe' aside) in ordinary type.  All four
% spans below were settled at 600 dpi, not from spacing.py alone: the two RUNs
% it reports are real and are here extended to their true boundaries, and its
% `Tvivl:' (1.49), `be-' (1.24) and `lute:' (1.52) are justification gaps, read
% at the image and NOT marked.  `klart' (1.23) IS letterspaced and is marked.
% Neither text layer records any of this, and the Google layer additionally
% scrambles this page's line order.
\opage{26}i Luften, og den Undersøgelse, der her skal anstilles, gaaer kun ud
paa at prøve, hvorvidt Ventilationen er tilstrækkelig. Den opstillede Sætning
er en Grundsætning; men det er neppe for dristigt at spørge, hvorpaa den
hviler. Den maa vel sagtens hvile paa de almindelige Sætninger, \emph{at
overhovedet to absolut uensartede Ytringer kunne være Ytringer af det Samme}
--- Tro og Viden ere jo principielle Ytringer ɔ: Ytringer af et Væsen, ---
\emph{og at overhovedet to absolut uensartede Principer kunne} --- skulle vi
springe „uden Modsigelse“ over? --- \emph{forenes i een og samme Bevidsthed}.
Lad os et Øieblik antage at denne sidste Sætning var sand, hvorfor staaer der
da „kunne“ og ikke „maae“? Kunne de forenes, saa kunne de ogsaa ikke forenes,
der er da ingen Nødvendighed, som tvinger. Og gjælder ikke det Samme, naar man
tager de to bestemte Principer: Tro og Viden? Uden Tvivl: altsaa man kan
forene dem og man kan lade være ɔ: Foreningen har ingen Almengyldighed.
% `Almengyldighed' with a single e here, against `Almeengyldighed' on p. 28.
% Both readings verified at the image; the book's own inconsistency, not
% normalised in either direction.
Men naar saa er, hvad bevirker da, at den, som forener dem, virkelig forener,
hvad der kan forenes? Her spørges, om dette lader sig gjøre \emph{klart}? Idet
Striden og Forsoningen mellem Tro og Viden af Prof.\ Nielsen er ført fra
Videnskabens Omraade over paa „Existensens“, er det blevet vanskeligt at have
Øie med, hvorledes den egentlig tænkes; thi paa Existensen kan man anvende det
gamle Udtryk, som Hegel brugte om Schellings Absolute: den er den Nat, hvori
alle Katte ere graae.
% --- p. 27 ---
\opage{27}Det saakaldte Existentielle, der gjør saa stærkt et Krav paa at
gjælde for Livets sande fuldblodige Concretion, er i Virkeligheden lige det
Modsatte, en puur Abstraction, nemlig det, der bliver tilbage, naar man fra
Menneskeaandens Tilværelse in toto subtraherer alt det Theoretiske; det er da
lige saa langt fra, at det efter den brugelige philosophiske Jargon sensu
eminenti existerende Menneske fører et sandt og ægte Menneskeliv, som den
Hund, hvilken en Naturforsker experimenterende har berøvet en Del af Hjernen,
% `Del' and `helt' and `naturligvis' on this page are the modern forms, read at
% the image; see ORTHOGRAPHY.md §2.  `puur' above is archaic, also as printed.
kan siges at leve Livet helt og virkeligt. Hvad der foregaaer i „Existensen“
lader sig da hverken let eller sikkert controllere. Men da Prof.\ Nielsen jo
fremfor Alt vil Klarhed, saa gaaer det naturligvis ikke an, at han paa noget
Punkt og mindst paa det afgjørende slukker Lyset, saa at Ingen kan see, hvad
der foregaaer, og først tænder det igjen, naar der er bragt ind og ud, hvad
der skulde bringes ind og ud i Theoriens Interesse. Hvad er det altsaa, som
bevirker den mulige Forenings Virkeliggjørelse? Ja hvad Andet er her at tye
til end den krasseste Umiddelbarhed, den personlige Trang, der idetmindste
\emph{kan} opfattes som en Sag, der er aldeles individuel. Dette er altsaa det
af saa stort et Chorus hilsede Philosophiens Mageløse: en høiere Brutalisme
% The colon after `Mageløse' is on the page at 600 dpi and in the KB layer; the
% Google layer drops it.
paa Basis af Instinct, umiddelbar Forvisning eller Andet af den Art. Skulle vi
antage denne Lære, lader os da først opgive at være Aander!

Idet man tager sin Tilflugt til den umiddelbare
% --- p. 28 ---
\opage{28}Praxis, der som saadan er uden Almeengyldighed, er man berettiget
til at tilfredsstille sin personlige Trang, som man lyster. Men skjøndt
Consequensen uafviseligt er denne, saa drager Prof.\ Nielsen den dog ikke
ligefrem. Hvilken er da den Kilde, hvoraf Philosophen øser det Indhold,
hvormed han udfylder den ene af de i den almindelige Sætning opstillede tomme
Rammer? Det er Aabenbaringen, Skriften, Religionen, hvormed da uden videre
menes den christelige Aabenbaring, Skrift, Religion, som denne opfattes i de
protestantiske Lande og navnlig i den seneste Tid er bleven opfattet i
Danmark. Dette kan den Videnskabsmand, der tager Prof.\ Nielsens Theori til
Rettesnor, kun gjøre derved, at han efter at have gereret sig som en subtil
Aand, i en Haandevending forvandler sig til en Kulsvier, der med et bredt Smil
svarer paa Ens Indvendinger, at hvad man troer, det troer man og ingen Snak om
den Ting --- for saa atter, naar den Pust er overstaaet, at kaste Vamsen,
vadske sig, rede sit Haar, tage sort Kjole paa og atter være subtil Aand. Hvis
En skulde begaae den Blasphemi at spørge, hvad \emph{Begreb} man skal gjøre
sig om en Aabenbaring i Almindelighed og den christelige i Særdeleshed, saa er
man Kulsvier og siger: Hvad der staaer, det staaer der. Staaer der, at man
selv skal bage sit Brød, saa skal man selv bage sit Brød, staaer der ikke
Noget om, at man selv skal brygge sit Øl, saa behøver man ingen Skrupel at
gjøre sig med Hensyn til Øllet. Staaer der, at man skal angre sine
% --- p. 29 ---
% CAPITAL Æ.  `Ære være ham' below was read at 600 dpi on this page: the book
% prints the Æ ligature.  The Google layer renders it `Ere' (ORTHOGRAPHY.md
% §3a); KB has it right.  Every other word on pp. 25--35 beginning with a
% capital vowel was checked against KB for the same failure and none was found.
% spacing.py reports no letterspacing here, and KB's `ved d e n ,' is NOT
% emphasis: read at 900 dpi, the letters of `den' are set at the same advance
% as those of `kunde' beside it, and the gap is the justification space before
% the comma.  Not marked.
\opage{29}Synder, saa skal man angre sine Synder; staaer der at Gud har skabt
Verden i 6 Dage, saa har han gjort det i 6 Dage, for stod der, at han havde
gjort det i 7 eller 14 Dage, saa havde han gjort det i 7 eller 14 Dage o.\ s.\
v. (Se Dr.\ Heegaards Fremstilling af Prof.\ Nielsens Lære „Indledn.\ til
Ethiken“ p.\ 392). Aabenbaringen er et Paradox, et absolut Paradox og som
saadant uden al Rationalitet. Spørg altsaa ikke: Hvad er en Aabenbaring?
eller: maa man drage Consequenser af det, som i en Aabenbaring er givet?
eller: hvor vidt maa man drage dem? Saadanne Spørgsmaal gjøre fornuftige
Mennesker, men det er ingen Konst at være fornuftig. Konsten er i samme
Øieblik, som slige Spørgsmaal gjøres, at sætte Ønskehatten paa Hovedet og saa
er man Kulsvier, hvis man ønsker det. Thi det forstaaer sig: er man berettiget
til at gaae endog kun en Linie udenfor de Grændser, Aabenbaringens Ord (6 ikke
7 eller 14) have draget, saa forsvinder det Overordentlige og saa er
Standpunktet forladt. --- Altsaa: man er Kulsvier, det er den ene Side af
% The colon after `Altsaa' is on the page and in KB; the Google layer drops it.
Sagen. Blev man staaende ved den, da kunde man endnu blive Gjenstand for
Sympathie og Agtelse. Respect for Kulsvieren, der i Sandhed er, hvad han
kalder sig! Ære være ham, naar han som Kulsvierne i Plougs „Fredericiaslaget“
gaaer løs paa sit Maal og slaaer for sin Fane „med Kulsviernæver og
Kulsviertro“! Men naar han paa samme Tid som han er Kulsvier, skal være den
videnskabeligt dannede Mand, da maa hans Fordring paa
% --- p. 30 ---
% FOOTNOTE, printed mark *).  It begins on p. 30 and RUNS OVER onto p. 31,
% where the continuation is set below the rule with no repeated mark; it is set
% here as one note at its mark on p. 30.  \endgraf would be needed for a second
% paragraph, but the note is a single paragraph.
%
% PRINTER'S DEFECT inside this note, transcribed as printed and NOT corrected:
% the quotation from Prof. N. is CLOSED with “ after `Brøde' but was never
% OPENED — the compositor set only a colon after `saaledes'.  Verified at 600
% dpi on the Google copy.  This is the batch's only quote-balance drift: the
% fragment therefore runs one „ short of its “, on purpose.
\opage{30}at betragtes som en helstøbt Skikkelse opbringe ethvert
sandhedskjærligt Sind.

Der staaer jo skrevet, at Døden kom ind i Verden ved Synden. Den orthodoxe
Opfattelse heraf er den, at Døden ikke staaer i noget væsenligt Forhold til
den menneskelige Organisme; havde de første Mennesker ikke syndet, vilde deres
Legemer have været uforgjængeligt, uopløseligt, uangribeligt af Sygdom og
% `Legemer have været uforgjængeligt' — the singular adjective against the
% plural noun is the book's; both layers and the image agree.  As printed.
Svækkelse. Men de syndede, og dermed var denne Fuldkommenhedstilstand
forspildt. Den moderne Videnskab er forlængst paa det Rene med, at Døden ikke
blot har rammet mangfoldige Dyreslægter, før Menneskene vare dannede, men at
Døden staaer i et nødvendigt Forhold til det organiske Liv. Her sige altsaa
Videnskaben og Orthodoxien hinanden skarpt imod; mod Orthodoxiens: „Døden er
kommen ind i Verden med Synden“, stiller Videnskaben sit „Nei! Døden er ikke
kommen ind i Verden med Synden“ paa samme Vis som den mod Orthodoxiens:
„Verden er skabt i 6 Døgn“ stiller sit „Nei! Verden er ikke skabt i 6 Døgn,
eftersom der er gaaet Tusinder af Aar, inden Jorden er bleven, som den
er“\fnmark{*)}{Naar Prof.\ N.\ (Grundid.\ s.\ Logik II p.\ 456) udtaler sig
saaledes: Det staaer fast, naar man blot ikke vil lade Videnskabens
„Forløsning“ gjøgle med Troen, at der har været Smerte og Død til i Millioner
Aar, før Mennesket fremtraadte paa Jorden. Og det staaer fast, naar man blot
ikke vil lade Troens „Forløsning“ gjøgle med Viden, at Døden er Syndens Sold
og Naturens Forkrænkelighed en Aandens Brøde“, da er Orthodoxiens haarde
Sætning her enten forflygtiget til en puur Aandrighed, til en dristig
Betegnelse af def uomtvistelige Slægtskab mellem Synd og Død, eller ogsaa er
% PRINTER'S ERROR, transcribed as printed: `def' for `det'.  Read at 900 dpi —
% the third letter carries the hooked top and two-sided crossbar of f, against
% the flat-topped t of `det' elsewhere in the same line of the note.  Both text
% layers independently read `def'.  The error is in the forme.
Forsoningen aldeles ufattelig; og naar Dr.\ Heegaard (Indledn.\ til Eth.\ p.\
394) erklærer, at man maa holde paa de 6 Dage, fordi man dog umuligt kan være
troende Geolog eller geologisk Troende, saa er dette et daarligt og
intetsigende Ordspil; thi den som --- det være sig nu som troende eller vidende
--- antager Nogetsomhelst om Jordens Tilblivelse, han har en Antagelse, som
Geologien kan corrigere.}. Nu er
% --- p. 31 ---
% The footnote block at the foot of this page is the RUN-OVER of the p. 30
% note, not a note of its own: it carries no mark.  It is set above, at its
% mark.  Three letterspaced `over' here, all confirmed at 600 dpi and all
% flagged by spacing.py (1.27, 1.30, 1.27).  The fourth `over', in `Nemesis
% over den' below, is ordinarily set and is NOT marked.
\opage{31}Prof.\ Nielsens Paastand jo den, at man uden Modsigelse kan forene
saadanne to Sætninger i sin Bevidsthed. Hvad er det da, som hidfører dette
forbausende Resultat? Svar: De ere absolut uensartede.

Man veed næsten ikke, hvorover man ved dette Svar skal undre sig mest,
\emph{over} at saadanne Sætninger virkelig skulle udgives for absolut
uensartede, eller \emph{over} at overhovedet absolut uensartede Phænomener
skulle kunne rummes af eet Jeg, eller endelig \emph{over}, at selve den
absolute Uensartethed skal være Grunden til det fredelige Forhold dem imellem.
Man fatter vel saa Meget, at hine Sætninger forsonligt skulle kunne bestaae
ved Siden af hinanden, fordi de saa omhyggeligt skulle holdes ude fra
hinanden, at det aldrig kommer til gjensidig Berøring, og man seer det maaskee
--- i Parenthes bemærket --- som en Nemesis over den i sin dialektiserende
% Both em-dashes here are on the page; the Google layer loses them entirely
% (it reads a stray `C' for the first and nothing for the second).
Bestræbelse forhen saa yderligt gaaende Dialektiker, at det her er en
Kriticisme, der skal frelse Theorien. Man begriber
% --- p. 32 ---
% FOOTNOTE, printed mark *).  Set at its mark, which the book places after
% `o. s. v.' and before the closing full stop.
\opage{32}vel ogsaa nok, at Sætningernes absolute Uensartethed kan være
Grunden til, at de ikke kunne modsige hinanden, da Betingelsen herfor jo
maatte være den, at de kunde sammenlignes, hvilket forudsætter Ensartethed
indtil en vis Grad. (Ensartethed indeslutter naturligvis Uensartethed og er
altid relativ, ellers var den Enshed.) Men hvad man aldeles ikke indseer, er,
at den absolute Uensartethed kan være Grunden til, at to Principer, Ytringer,
Sætninger kunne forsones; tvertimod er den en afgjørende Grund til, at de
aldeles ikke kunne forsones, da de ikke engang kunne komme i Forhold til
hinanden. Absolut Uensartethed er et andet Udtryk for absolut Uforenelighed.
Prof.\ Nielsen kunde derfor lige saa godt have givet sin Sætning denne Form:
Absolut uforenelige Principer kunne forenes eller være ikke absolut
uforenelige --- Alt uden Modsigelse; han har sagt det, men med andre Ord. Da
% The em-dash before `Alt' is on the page and in KB; the Google layer drops it.
jeg i sin Tid i „Fædrelandet“ stemplede det som en Selvmodsigelse at lade Tro
og Viden, naar de bestemtes som absolut uensartede Principer, forenes i samme
Bevidsthed, reiste man sig fra flere Sider redebon til med tillidsfuld
Overlegenhed at vise mig, at mit Angreb paa Prof.\ Nielsen beroede paa en
fuldkommen Misforstaaelse; hvor jeg satte et \emph{skjøndt}, der satte Prof.\
et \emph{fordi} o.\ s.\ v.\fnmark{*)}{se Høffding: „Philosophie og Theologie“
p.\ 41. R.\ Schmidt: „Om Selvmodsigelsen o.\ s.\ v.“ p.\ 17. Heegaard:
Indledn.\ t.\ d.\ r.\ Ethik p.\ 389.}. Der gives en Slutningsmaade, der er
% `p. 17' is lower case on the page; the Google layer capitalises it (`P. 17').
% Both \emph{} spans here were read at 600 dpi.  spacing.py caught only
% `fordi' (1.29); `skjøndt' is letterspaced too, and KB's `s k j ø n d t'
% corroborates it (ORTHOGRAPHY.md §3g).
ligesaa farlig og mindst lige%
% --- p. 33 ---
% The page turn falls INSIDE `ligesaa'.  The lone numeral at the foot of this
% page is printer's signature 3, not text, and is not transcribed.
% No letterspacing on pp. 33 or 34.
\opage{33}saa almindelig som Fusentastens: jeg kan ikke forstaae den Mand,
følgelig har han Uret, det er Eftersnakkerens: jeg forstaaer denne Mand,
følgelig har han Ret. Lærlingen begynder med at staae undrende, usikker, fuld
af Spørgsmaal og Tvivl overfor Lærerens Tale, men naar det saa --- ofte
pludseligt --- gaaer op for hans Bevidsthed, hvorledes Læreren egentlig tænker
sig Sagen, paa hvilken Maade han vil have det Udviklede forstaaet, saa bliver
han tidt og ofte i sin Glæde over at fatte, hvad der for Mange og saa nylig
for ham selv var en Gaade, i samme Øieblik en Tilhænger og des mere ubetinget,
jo mere uselvstændig han af Naturen er. En saadan Lærling udleder naturligvis
enhver Indsigelse mod hans Mesters Lære i det heldigste Tilfælde af
Misforstaaelse og Uvidenhed, i værste Fald af en moralsk Depravation, og man
maa holde hans Standpunkt Noget tilgode: at En kunde forstaae, hvad Læreren
vilde, og endda have Noget at indvende derimod, det er ikke faldet ham ind i
hans dristigste Drømme. Jeg skulde da altsaa ikke have forstaaet, at den
absolute Uensartethed netop var Betingelsen for den meget omtalte Forsoning.
Hvorvel jeg allerede dengang udviklede min Anskuelse temmelig omstændelig, og
navnlig ikke som det for nylig usandfærdigt er bleven sagt af Hr.\ Clemens
Petersen trak mig tilbage, da Modparten „gjorde Mine til at føre Sagen ind i
en gjennemgaaende Forhandling“, vil jeg dog her optage dette Stridspunkt
paany. Sagen er ganske simpel. Hvis saadanne to Sætninger,
% --- p. 34 ---
\opage{34}som a) „Verden er skabt i 6 Døgn“, b) „Verden er ikke skabt i 6
Døgn“ ere absolut uensartede, saa er dette det Samme som, at den ene er en
Ytring af et Væsen (essentia), der er absolut uensartet med det Væsen, hvoraf
den anden er en Ytring. Fiat applicatio --- den R.\ Nielsen, der udsiger
% `Fiat applicatio', `essentia', `in toto' and `sensu eminenti' are set in the
% ordinary roman with no contrasting face; no \textit{} — see the header.
Sætningen a, er en absolut anden end (uensartet med) den, der udsiger
Sætningen b; at kalde dem begge med samme Navn er en uskyldig Morskab ligesom
Hostrups Indfald i „Soldaterløier“ at kalde alle tre Arvinger Peter. At Prof.\
Nielsen kalder sig selv med samme Navn i begge Tilfælde, er en Caprice, man
kan lade ham, naar man blot passer, ikke at lade sig vildlede deraf. Her er
imidlertid ikke blot to Personer, men to Personer, der høre hver til sin med
den anden absolut uensartede Verden; thi enhver Spaltning af et Subject er en
Spaltning af hele den objective Verden, der speiler sig i Subjectet, eftersom
dette dog i sidste Instans er en Ytring af hin; er Subjectet delt, saa kommer
det principielt deraf, at den hele Verden har spaltet sig (essentia) forud.
Man bør da ikke spørge, hvorledes det staaer sig med en Mands Selvbevidsthed,
naar han huser saadanne hinanden modsigende Grundanskuelser; thi en saadan
Mand kan ikke existere. Den absolute Uensartethed er den absolute Dualisme;
men den absolute Dualisme er en Maskeradespøg, en Attrappe, der skjuler en
virkelig Ensartethed og Enhed.

Da man fra det engang indtagne Standpunkt ikke gaaer udenfor de Træk,
Aabenbaringen indbefatter,
% --- p. 35 ---
% FOOTNOTE, printed mark *).  The `3*' at the foot is printer's signature, not
% text.  `udenfor' and `paa' are letterspaced (spacing.py 1.26 and 1.26, both
% confirmed at 600 dpi); spacing.py's `Eventyr,' (1.42) is the justification
% gap before the comma and is NOT marked.
\opage{35}kan man ikke yderligere vise dets Urimelighed: den christelige
Aabenbaring indeholder jo nemlig en --- naar Hensyn tages til den nye Lære ---
forbausende Mængde af fornuftige Sandheder, der ikke lade sig studse op til de
nu saa yndede Paradoxer\fnmark{*)}{Angaaende dette Punkt er der i vor nyere
Litteratur bleven drevet den uheldigste Misbrug med Ord, idet enhver i det nye
% A small raised mark stands between `nyere' and `Litteratur' on the page and
% is recorded by BOTH text layers as an apostrophe.  Read at 900 dpi it is a
% speck of ink at ascender height, not a sort: Danish takes no apostrophe here.
% Not transcribed.  (Compare ORTHOGRAPHY.md §3e, where the Google layer alone
% promoted a speck to a full stop; here both engines did.)
Testamente udtalt Sandhed betegnes som ikke opkommet i noget Menneskes Hjerte
eller Hjerne. Enhver Sætning, som Mennesker kunne opfatte og forstaae, maa
Mennesker imidlertid fra først af have undfanget; thi at frembringe og opfatte
er een og samme Sag seet fra to forskjellige Sider. Dernæst betegnes enhver
saadan Sandhed som paradox, det Humane modsat: saaledes skal det f.\ Ex.\ være
en paradox, en positiv-christelig Sætning den, at man skal elske sine Fjender.
Men i Virkeligheden seet er det slet ikke mod Forstand eller Fornuft, men mod
Sandseligheden, den strider; at man ikke skal gjengjælde Ondt med Ondt er et
rent Fornuftbud. Var Alt det paradoxt, som kommer i Strid med vor sandselige
Drift, da var Fornuften selv det Paradoxe.}. Derimod er der, naar man, hvad
man er berettiget til at gjøre, stiller sig \emph{udenfor} Standpunktet
(christelig Aabenbaring), men \emph{paa} dets Grundsætning i dennes
% `Grundsætning' is set with a visible gap after `Grundsæt' — a loose space in
% the forme, recorded by both layers.  It is one word and is set as one word.
almindelige Form (Muligheden af absolut uensartede Ytringers Forening) et
uendeligt Spillerum for de vildeste Forsøg paa at tilfredsstille den praktiske
% The Google layer invents a full stop after `vildeste'; the page has none.
Interesse, den personlige Trang. Saadanne ere allerede skete; se de
mangfoldige positive Religioner, der jo i Eventyrlighed overgaae de egentlige
Eventyr, og dog netop
% --- p. 36 ---
% Continues from p. 35, which ends `og dog netop' at a clean word boundary
% (the rest of p. 35 is taken up by its footnote and the signature `3*').
% CAPITAL Æ: `Æren' read at 900 dpi on the Google image (the ligature is
% unambiguous).  The Google layer reads `Eren' here, as the audit predicted;
% KB reads `Æ ren'.  Every other word-initial capital vowel in pp. 36--46 was
% checked at the image as well (Ens, Ende, Exempel, Element, Eftertrykket,
% Efterskrift, Eng, En, Egnens, End, Et, Ensartethed, Er, Ex.) — all are plain
% E/A, no further Æ in this range.
% `lendset' is as printed (l-e-n-d-s-e-t, read at 900 dpi; both layers agree).
% No letterspacing on this page: spacing.py detected none, and KB's `jo b ed re'
% is ordinary ABBYY splitting — the phrase is normally set on the image.
\opage{36}ere Frugten af den dybeste personlige Trang. Alle disse maa Prof.\
Nielsen lade gjælde; om de have et fornuftigt Ideeindhold eller ei, gjør jo
fra et overfornuftigt Standpunkt Intet til Sagen. Man kan ikke engang sige:
% The colon after `sige' is on the image and in the KB layer; the Google layer
% drops it (cf. the dropped colon at p. 59 recorded in ORTHOGRAPHY.md §3e).
jo galere, jo bedre; thi Alt kan fordre lige Rang, hvor den personlige Trang,
renset eller lendset for al Fornuft, er Maalestok. Det, at man som Dyrker af
sin Gud og Dyrker af sin Videnskab paastaaer absolut stridige Udsagn, er fra
Midten af det 19de Aarhundrede at regne ikke til Hinder for, at man udtaler
den klare, ja den høieste Sandhed. Det er Danmark, som har Æren af denne
Opdagelse, ved Siden af hvilken man maaskee kan stille N.\ F.\ S.\
Grundtvigs, men visselig hverken Ole Rømers eller H.\ C.\ Ørsteds.

Men man behøver ikke at blive staaende herved, de positive Religioner have
ikke udelukkende Ret til Tilfredsstillelsen af den antirationelle Trang. Det
er min Trang og ingen Andens, jeg skal fyldestgjøre; om jeg derved kommer til
at leve i et Mylr af absolut uensartede Ytringer, hvad gjør det? Det strider
mod Fornuft, mod Tanke, Fornuft! Tanke! De have ingen Stemme, hvor det
% `Tanke,' — COMMA, not a full stop.  The two layers disagree (Google comma,
% KB period) and the glyph is a filled blob, so it was measured at 900 dpi:
% 19x27 px, identical to the certain comma after `Verdensprincip' on p. 37,
% whereas this page's periods measure 19x19 and its open commas 21x40.  The
% counter has simply filled with ink.  KB is wrong here.
gjælder det Høieste. Man kan vælge hvilken modsigende Antagelse, man vil;
naar man erklærer den fordret af Ens personlige, Ens ethiske Trang, kan Prof.\
Nielsen Intet indvende mod den. Til Anskuelighed heraf blot et enkelt
overfornuftigt Træk. Hvis En har den personlige Tro, at Martin Luther var en
Incarnation --- vel ikke af det
% --- p. 37 ---
% No footnote, no letterspacing (spacing.py: none).
\opage{37}hele Verdensprincip, men af Siriusprincipet f.\ Ex., og hvis han saa
for den Samme som vidende har været den bekjendte Reformator, der aldeles
Intet har med Siriusprincipet at gjøre, hvad vil det da gavne at kalde Sligt
urimeligt og unaturligt? Sandsynlighed og Fornuft have paa dette Punkt jo
Intet at sige. Eller skulde Philosophen virkelig ville søge Beroligelse i den
% `Eller skulde' — no comma after `Eller'; KB inserts one, the image has none.
Tanke, at Ingen vil falde paa rene Urimeligheder, skulde han blive reduceret
til at appellere til Sandsynligheden, hin det Paradoxes Arvefjende,
Sandsynligheden af at Ingen vil falde paa Noget, der i den Grad strider mod
Fornuften? Thi denne er jo dog den eneste, der kunde have Noget imod saadanne
Antagelser at indvende, Overfornuften vilde jo see Alt, og see, det var saare
godt. Skulde Principernes Forfægter virkelig blive indskrænket til at ville
til det Yderste have i Principet, hvad han paa det Bestemteste vil have sig
% `Yderste' — CAPITAL Y, read at the image.  The Google layer lowercases it;
% KB has the capital.
frabedt i Virkeligheden?

Da den hele besynderlige Lære saa vanskeligt lader sig tydeliggjøre, naar man
gaaer lige løs paa Sagen, have Tid efter anden Tilhængere af Prof.\ Nielsen
vist Publicum den Velvillie at oplyse Forholdet ved Billeder eller Analogier.
Uden at opholde os ved, hvor mislig en slig Oplysning efter sit Væsen er,
ville vi blot see, med hvad Held, den er foretaget. For ikke at lade noget
Herhenhørende uomtalt, vil jeg minde om, at Hr.\ Clemens Petersen engang for
nogle Aar siden fremstillede Forholdet mellem Forestilling og Begreb paa det
religiøse Om%
% --- p. 38 ---
% `Del' (not `Deel') as printed — see ORTHOGRAPHY.md §2; read at the image.
% LETTERSPACING inside the Heegaard quotation: `kan', `fordi', `uagtet', all
% three confirmed at 300 dpi (wide, even inter-letter gaps) and flagged by
% spacing.py at 1.25/1.36/1.32.  spacing.py also flagged `(det' at 1.25; the
% image shows `(det contemplative)' ordinarily set — the parenthesis inflates
% the predicted width.  REJECTED.
% No footnote on this page.
\opage{38}raade ved følgende Billede. Naar, hed det sig omtrent, En seer en
Møllevinge dreie sig om sin Axe, og man spørger den Paagjældende, om den
yderste eller inderste Del af Vingen har den hurtigste Bevægelse, da vil han,
hvis han retter sig efter det sandselige Indtryk, uvilkaarligt svare: de gaae
lige hurtigt [?]; begrebsmæssigt opfattet gaaer derimod naturligvis den
% `[?]' is printed, square brackets and all — Brandes's interpolation into
% Clemens Petersen's parable, not an editorial mark.
yderste Ende meget hurtigere end den inderste. Saa lidet, hed det nu, som man
burde „pine Forestillingen til at see Begrebet“, saa lidet burde man
paatvinge Tanken Forestillingens Illusion: Saaledes ogsaa i det Religiøse;
der gives een Sandhed for Forestillingen, en anden, hin første modsigende,
for Tanken, ingen af disse gjør den anden noget Afbræk. --- Dette var
Lignelsen; man seer, den svarer til Prof.\ Nielsens Theorie, og at Hr.\
Petersen altsaa, hvad ikke Mange vide, har anden Fortjeneste af denne Sag end
den at være Professorens Panegyrist. Tør man slutte fra hans Lignelse, da
blive de religiøse Forestillinger lutter Sandsebedrag af tvivlsom
Uundgaaelighed. Hr.\ Petersen synes imidlertid at have villet godtgjøre noget
ganske Andet.

Dernæst har Dr.\ Heegaard (Indledn.\ til Ethiken p.\ 389) leveret en
oplysende Parallel. Hans Ord ere disse: „Paa dette (det contemplative) Trin i
% `Parallel' — the r is damaged in this copy (Google reads `Parallel', KB
% `Paiallel'); the image shows r with a broken leg.  Set at p. 39 line 1 of
% the Indhold-style spelling `Parallel' again, undamaged, which settles it.
Bevidstheden \emph{kan} der ingen Strid indtræde mellem Tro og Viden, netop
\emph{fordi} de ere uensartede og ikke, som man ved en besynderlig
Misforstaaelse oftere (sic!) har villet gjøre gjældende, \emph{uagtet} de ere
uensartede
% --- p. 39 ---
% LETTERSPACING: `kunne', confirmed at the image (spacing.py 1.23).
% spacing.py also flagged `deten' at 1.25 — a mis-segmentation, no such word;
% nothing else on the page is letterspaced.  REJECTED.
% No footnote on this page.
\opage{39}Principer; dette „uagtet“ involverer nemlig, at Troens og Videns
Sphærer skulde \emph{kunne} begrændse hinanden, hvilket er umuligt, saalænge
Bevidstheden er oversvævende. Vi vælge et Exempel: Den contemplative
Bevidsthed kan uden Modsigelse rumme paa eengang et æsthetisk og et
astronomisk Problem, ikke uagtet, men netop fordi de ere uensartede; den kan
naturligvis ikke til samme Tid have dem begge i Bevidstheden, men den kan
fordybe sig i ethvert af dem, saameget den vil, uden at det ene Problem
nogensinde kan lide Afbræk ved det andet, just fordi de vise hen til
uensartede Udgangspunkter.“ Det turde være unødigt at give nogen Commentar
til dette Sted, det maa vel være aabenbart selv for den lettest Bedaarede, at
der her end ikke er Spor af Analogi tilstede. Man vil iøvrigt bemærke, at
Dr.\ Heegaard omhyggeligt undlader at kalde Uensartetheden absolut. Da jeg
først bestred Prof.\ Nielsens Brug af dette Udtryk, kappedes DHrr.\ Schmidt
og Høffding om at forklare mig, at det just var deri, Hemmeligheden stak;
havde Uensartetheden været relativ, da var det en ganske anden Sag, da vilde
Forsoningen have været umulig. Dr.\ Heegaard synes dog at have opdaget, at
man ved at lægge Eftertrykket paa dette Punkt kom fra Scylla i Charybdis.

Den tredie og mest omtalte Parallel er den, Hr.\ Rud.\ Schmidt har draget
mellem dette Forhold, og Forholdet mellem Poesiens og Virkelighedens Sandhed.
Denne Analogi har noget større Interesse, da
% --- p. 40 ---
% `helt og holdent' — the modern form, as printed (ORTHOGRAPHY.md §2 records
% `helt' at p. 40); read at the image.
% No letterspacing (spacing.py: none).  KB's `N ei', `U d try k' and
% `U dtryk' are ABBYY splitting, not Sperrsatz: all three are ordinarily set
% on the image.  No footnote on this page.
\opage{40}Kierkegaard maaskee flygtigt har antydet den, se hans Benyttelse af
% `se', single e, as printed — contrast `See' opening the p. 42 footnote and
% `See vi da nu tilbage' on p. 43, both double-e.  All three read at the image.
Gorgias's: Den Bedragne er visere end den Ikke-Bedragne, smlgn.\ Afsluttende
Efterskrift p.\ 349. Spørgsmaalet er da først og fremmest, om de tvende Led
ere absolut uensartede. Man maa strax svare Nei; thi Benegtelsen ligger
allerede i selve de brugte Udtryk: have begge Sandhed, saa ere de begge
Momenter i Sandhedens Rige, og det er da saa langt fra, at der her kan være
Tale om absolut Uensartethed, at der endog maa være et Element, hvori Poesie
og Virkelighed ere absolut ensartede. Hvis nemlig Sandhed er et Begreb, der
virkelig kan anvendes paa dem, eller, for at bruge et inderligere Udtryk, en
Væsensbestemmelse ved dem, saa maa denne helt og holdent gjennemtrænge dem
begge, gjennemtrænge enhver af deres andre Bestemmelser. Begge ere da
Ytringer af en dem omfattende Totalitet. Saavidt i Almindelighed. Men det
viser sig ikke mindre klart, naar man griber dette Forhold i et bestemt
Exempel. Hr.\ R.\ Schmidt har brugt et saadant, som Prof.\ Nielsen ved ikke
at underkjende det har lyst i Kuld og Kjøn. Han har anført den Folketro, at
naar Engen damper en Sommermorgen, da er det Mosekonen, der brygger: ligesom
der i dette Tilfælde for den Dannede ingen Modsigelse er mellem den poetiske
Tro og den naturvidenskabelige Erkjendelse, saaledes er der heller ingen
Modsigelse mellem Troens og Videns uensartede Antagelser. Men for det Første
er Phantasiengen og den virke%
% --- p. 41 ---
% LETTERSPACING: `jeg' and `Staten', the contrasted pair in the Fuursøen
% example — confirmed at the image (spacing.py 1.24 / 1.30).
% Both em-dashes in `Mosekonen --- absolut naturligvis ---, hvad' are on the
% image; the Google layer renders them as hyphens and misreads the comma after
% the second as a quotation mark (ORTHOGRAPHY.md §3d records this very line).
% `naturligvis', single i, as printed.  No footnote on this page.
% NB the Google layer opens a spurious blank line before `Menneskes
% Opløftelse'; the image shows the text running on without a paragraph break.
\opage{41}lige Eng ikke absolut uensartede og følgelig ligesaa lidt de
tilsvarende Opfattelser. Skil Mosen fra Mosekonen --- absolut naturligvis
---, hvad er saa den sidste for Indbildningskraften? Folketroens
Personification er jo ikke Andet end Virkelighedens Reflex i Phantasien, der
uophørlig maa arbeide med det Virkelige for at frembringe Noget, der ikke er
Virkelighed og dog er Virkelighed; for at forstaae dette, maa man bruge
Dialektik, men Dialektiken er jo i denne Sag een Gang for alle devoveret. For
% `devoveret.' ends the sentence with a full stop on the image; KB reads a
% comma.
det Andet er her ingen Analogi; thi, da den Eng, som Folkepoesien har, i
Almindelighed vil kunne betegnes som en Phantasieng, men den, som den
naturvidenskabelige Opfattelse rettes paa, er den virkelige Eng, forholde de
sig ikke til een og samme Gjenstand. Man kunde da ligesaa godt kalde det en
Modsigelse af samme Art som den mellem Tro og Viden, naar jeg paa een Gang
siger, at \emph{jeg} eier Fuursøen (Maleri af N.\ N.) og at \emph{Staten}
eier Fuursøen (den virkelige). En anden Sag er det, ifald man paastaaer, at
den virkelige Engs Damp er en Virkning af, at Mosekonen brygger, og saa
tillige siger, at den lige saa fuldt er en Virkning af physiske Aarsager: Da
have vi en Analogi, en Modsigelse som den, hvorom Talen er; thi naar f.\ Ex.\
Orthodoxien antager et Menneskes Opløftelse fra Jorden og gradvise Svæven ind
i den dogmatiske Himmel, og Videnskaben benegter dette Under først i Kraft af
Tyngdeloven, dernæst fordi man, om man end løftes nok saa høit i
% --- p. 42 ---
% FOOTNOTE on this page, printed mark *) — one note only.
% `dels' twice, the modern form, as printed (ORTHOGRAPHY.md §2); read at the
% image.  `nødvendigvis', single i.  `Stof' with a single f, against `Stoffet'
% two lines later; both as printed.
% No letterspacing in the body.  KB splits nearly every word of the footnote
% (`In d b ild n in g sk ra ft'), which is ABBYY degrading in the small type,
% not Sperrsatz — the footnote is ordinarily set on the image.
\opage{42}Veiret, dog aldrig kommer ud af denne Verdens astronomiske Himmel
og ind i den transscendente, saa er det det selvsamme Punkt, den selvsamme
Himmelfart, hvorom de begge tale. --- Den søgte Analogi findes da heller ikke
% PRINTER'S DEFECT, transcribed as an em-dash: the dash after `tale.' took ink
% only along its lower edge and prints as a short low stroke at the baseline,
% about a third the length of the two well-inked em-dashes lower on this same
% page.  Read at 900 dpi.  The Google layer records a hyphen, KB drops it
% entirely.  The construction (full stop + em-dash + capital) is the book's
% own, cf. `noget Afbræk. --- Dette' on p. 38.
her. Men uanseet dette, selv naar man vil nøies med i Analogiens Interesse at
betone den Lethed, hvormed i det brugte Exempel to uensartede Opfattelser
glide hen over hinanden uden gjensidig Paavirkning, da glemmer man dels ---
og dette er meget slemt --- at Virkelighedserkjendelsens Fremskriden
umuliggjør Brugen af visse tilforn vel brugelige
Phantasiudtryk\fnmark{*)}{See Ørsted: Aanden i Naturen II.\ p.\ 4. „For en
Indbildningskraft, der har tilegnet sig et levende nærværende Billede af
Verdenssystemet, er Udtrykket Jordens Grundvold ikke bedre end Grundvolden
for en velophængt Lysekrone, men om muligt endnu mindre passende“.}, dels
overseer man, at hin saakaldte Gliden hen over hinanden i Virkeligheden ikke
er noget Saadant, men snarere maatte bestemmes som den ene Opfattelses
Indtrængen i den andens Lacuner eller Porer, muliggjort ved, at ingen af dem
er ret fuldstændig. Derfor har f.\ Ex.\ Digteren des friere Haand over det
historiske Stof, jo mere ubekjendt det er for hans Publicum; hvis dette maa
forudsættes bekjendt med hvert lille Træk i Stoffet, da vil han nødvendigvis
være meget bunden. Det vilde imidlertid være, hvad Kierkegaard kalder
„Parenthesdurchløberi“ at indlade sig paa en nøiere Drøftelse heraf i
% --- p. 43 ---
% `momentvis', single i (ORTHOGRAPHY.md §1).  `duæ veritates' is set in the
% ordinary roman with no contrasting face, so it takes no markup.
% No letterspacing (spacing.py: none).  No footnote on this page.
\opage{43}Anledning af det Spørgsmaal, hvorpaa det kommer an. Det Afgjørende
er og bliver, at den poetiske Illusion som et fuldstændigt Unicum aldeles
ikke er nogen Analogi: den er vel at mærke kun poetisk, naar den momentvis
afvexlende er og er ophævet, den er vel at mærke kun poetisk som bevidst; thi
Fuglene, som hakke i de malte Bær, eller Lazaronerne i Neapels Smaatheatre,
% `hakke' on the image; KB reads `bakke'.
der med Fagter røbe Helten, hvor Skurken har skjult sig, ere ikke naaede til
den poetiske Illusion. Det Besynderligste er imidlertid, at Forsvarerne af
denne Analogi først vilde være rent galt farne, hvis man vilde holde den for
Alvor fast; thi den, der vil tillægge en eller anden bestemt religiøs
Forestilling den fuldeste Virkelighed, han har dog handlet sig selv imod,
saasnart han jævnfører den med det rent Imaginære --- det være nu nok saa
skjønt og nok saa betydningsfuldt, baade langt skjønnere og langt mindre
taaget end Mosekonens Bryg. Bliver Mosekonen derimod forvandlet til en
virkelig Person, da have vi en af de antirationelle Sandheder, for hvilke
Prof.\ Nielsen consequent maa bøie sig, og da have vi i fuld Gang Spillet med
de to, hvilke Prof.\ Nielsen forfægter paa godt Dansk, men bekjæmper, naar de
optræde som Middelalderens duæ veritates.

See vi da nu tilbage paa det, som her er indvendt mod den nye Lære, hvad
Betydning faaer saa den hele Opstilling af absolut uensartede Principer,
% --- p. 44 ---
% LETTERSPACING: `ere' in `naar absolut uensartede Gjenstande ere, saa ere de
% ikke absolut uensartede' — the first `ere' only, confirmed at the image
% (spacing.py 1.25).  The second is ordinarily set.
% `væsenligt' as printed.  No footnote on this page.
\opage{44}Ytringer, Sandheder eller hvad man skal kalde disse ilde stedte
Væsener? Skulde det ikke forholde sig med dem som med de Byer, hvilke
Potemkin viste sin Keiserinde som Vidnesbyrd om Egnens blomstrende Tilstand?
De vare nemlig i egentlig Forstand kun opstillede, thi de existerede ikke.
Subjectivt seet kan der for os umuligt gives absolut uensartede Gjenstande;
de maatte jo som Aarsager betragtede fremkalde absolut uensartede
Bevidsthedsytringer i det følgelig aldeles spaltede Subject. Kan da objectivt
seet nogensomhelst Ytring, der fremgaaer af et Tilværelsesprincip eller af et
Indbegreb af Objecter, der danne en Totalitet, være absolut uensartet med
nogensomhelst anden? Tænk det og Du kommer i Selvmodsigelse, idet Du sætter
en Helhed og nægter hvad der følger deraf. At være, det er at være indbefattet
i en Tilværelse; naar absolut uensartede Gjenstande \emph{ere}, saa ere de
ikke absolut uensartede. Man gjøre en Prøve med hvilkesomhelst Gjenstande,
Ytringer, Phantasier eller hvad man vil, og man vil finde, at selv om de
udelukke hinanden i en vis Henseende, gjøre de det dog i en anden Henseende
ikke, idetmindste forsaavidt som de begge kunne falde i eet Subject. Prøv:
„en gul Symphoni“, det kan ikke siges og dog ere Tone og Farve væsenligt
ensartede som Bølgephænomener; „en ligesidet retvinklet Trekant“, det gives
ikke og dog betegne begge Tillægsord mulige Bestemmelser ved Be%
% --- p. 45 ---
% FOOTNOTE on this page, printed mark *), keyed to `Trekant*)' in the first
% line.  `Del' (not `Deel') at `for en stor Del Prof. Nielsen selv', as
% printed (ORTHOGRAPHY.md §2); read at the image.
% No letterspacing (spacing.py: none); KB's spaced-out footnote is ABBYY
% degrading in the small type, as on p. 42.
\opage{45}grebet Trekant\fnmark{*)}{End ikke saadanne udprægede Modstandere
som Digterne Hr.\ Henrik Scharling og Hr.\ Rud.\ Schmidt ere absolut
uensartede. Der gives tvertimod Punkter, hvor en absolut Ensartethed finder
Sted imellem dem. Et saadant rammer Hr.\ Schmidt med megen Sikkerhed, naar
han i sin Bog: „Er Hr.\ Lic.\ Henrik Scharling en alvorlig Mand?“ pag.\ 80
fremsætter den Formodning, at Hr.\ Scharling nu „som sædvanlig har digtet os
noget Slet“.}. Tag det, som aldeles ingen Betydning har: det Meningsløse, det
% The colon after `har' is on the image and in the KB layer; the Google layer
% drops it.
Ikkeværende, skil det absolut fra det, som er, og det kan ikke engang
udsiges; selv dette, der synes at maatte have Privilegium paa at existere,
sondret fra sit Modsatte, er intet Andet end en Reflex af dette. Men naar
Uensartetheden i absolut Forstand selv her ikke har nogen Magt, hvor meget
mindre har den det da der, hvor Alt tillægges Tilværelse og Realitet.

Man vilde gjøre mig høilig Uret, hvis man antog, at jeg ved at fremsætte rent
elementære Sandheder som disse, noget Øieblik har indbildt mig at lære en
Tænker som Prof.\ Nielsen Noget, han ikke selv fuldt saa godt vidste, Noget,
han ikke selv har udtalt, længe før jeg for min Part havde Anelse derom. Jeg
skylder for en stor Del Prof.\ Nielsen selv, at jeg er istand til at modsige
ham; men det forholder sig jo ikke saaledes med Sandheden, at den, der som et
af dens Organer har meddelt den, senere efter Behag kan tage sin Gave
tilbage, naar den som Vaaben selv i den Svagestes Haand vilde være stærk nok
til at for%
% --- p. 46 ---
% Last page of chapter 2.  Eight lines of text, then a short centred double
% rule closing the chapter; p. 47 opens chapter 3 and is not part of this
% batch.  The page carries no folio-side ornament and no footnote, and nothing
% on it is letterspaced (spacing.py: none).
\opage{46}styrre en dristig anlagt Theorie, opført paa et skrøbeligt
Grundlag, en Theorie, der vil hæmme den aandelige Fremskriden, der vil spærre
Veien for enhver Stræben efter at bringe Harmoni tilveie mellem Fornuft og
Religion, mellem Menneskeaandens intellectuelle og dens ethisk-religiøse
Krav, og som saaledes forstyrrer enhver Arbeiden paa Løsningen af en af
Menneskeslægtens store Fremtidsopgaver.

\begin{center}
\rule{2cm}{0.4pt}\\[-4pt]
\rule{2cm}{0.4pt}
\end{center}

%% ============================================================
%%  3. DET ANTIRATIONELLE  (printed pp. 47--54 = PDF 51--58)
%%  Head as printed: `3.' / `Det Antirationelle.'
%%  Indhold reads `Det Antirationelle', without number and without stop.
%%  NB the Google text layer renders this head with a space before the stop
%%  (`Det Antirationelle .'); the image shows no such space.  OCR artefact.
%% ============================================================
\chapter*{3.\ Det Antirationelle.}
\addcontentsline{toc}{chapter}{3. Det Antirationelle}
\markboth{Det Antirationelle}{Det Antirationelle}
\label{ch:3}

% --- p. 47 ---
% Chapter opening: the head `3.' / `Det Antirationelle.' is already set in the
% skeleton and is NOT repeated here.  The page drops its folio and opens with
% an enlarged capital `M'; the enlarged capital is not marked up.
% NO LETTERSPACING on this page.  spacing.py flagged `Det' (1.26), `Man' (1.23)
% and `Liv' (1.60); all three were rejected at the image.  `Det' is the display
% head, `Man' measures wide only because of the enlarged M (the `an' is tight),
% and `Liv' is ordinary type — the detector caught a speck of foul ink above
% the `v'.  KB's `L itteratur' is ABBYY splitting, not Sperrsatz: confirmed
% solid at the image.
% `Bedste' carries a small ink blot below the B; the reading is unambiguous.
\opage{47}Man kan sige sig selv, hvad Følgen maa blive, naar en Tænker overseer
hine elementære Sandheder og til Bedste for en personlig Trang, der naar der
tillægges den Almeengyldighed, bliver en Særhed, der piner og trykker Andre,
forsøger paa at holde det Uholdbare fast. Følgen har allerede vist sig. Den er
et Begrebsanarchie, en Begrebsforvirring, der undergraver den Grund, hvorpaa
vor Videnskab bygges, og for hvilken Theorierne hos Dr.\ Heegaard sikkert er
det mest fortvivlede Udtryk, som nogen Litteratur i den Retning har at opvise.
Thi hvad have vi nu? Vi have den velbekjendte Metaphysik og den velbekjendte
Logik med deres Begreber, fribaarne, fødte af Natur og Fornuft, levende deres
Liv i en Verden, i hvilken de ere indordnede som Led og hvor deres Væsen
udfoldes og kommer til sin Ret.
% `Ret.' — FULL STOP, read at the image.  The layers disagree (Google period,
% KB comma); Google is right here.
Og saa have vi paa den anden Side en Verden, hvor de maa gjøre den laveste
Trælletjeneste, den Tjeneste nemlig, i hvilken Trællen maa fornægte sin
Herkomst og sit Væsen. Thi disse samme Begreber og Magter anvendes jo i den
antirationelle Verden saaledes, at de slet intet
% --- p. 48 ---
% ONE letterspaced word: the second `Intet' (`ja at de I n t e t skulle
% udtrykke'), read at the image and corroborated by KB's `de I n t e t'.  The
% FIRST `Intet' three words earlier (`at de Intet skulle') is set solid, and so
% is the whole of the following line — checked at the image, because Brandes
% normally letterspaces a clause, not a word.  Here he does not.
% THREE EM-DASHES on this page, all taken from the image: after `Consequens',
% and the pair round `som man siger, indtil Skjorten'.  The Google layer drops
% all three (it prints a blank or nothing); KB has them.
\opage{48}Rationelt maa udtrykke. Gjøre de nemlig dette, saa forvandle de med
det Samme de impotente Magter, for hvilke de skulle gjøre Hoveri, de saakaldte
Kategorier: det Paradoxe, det Antirationelle, det Positive, til rationale
Størrelser, hvad disse frygte for som visse Væsener ifølge Sagnet frygte Dagen.
Hine Lysets Børn tvinges da til at udslukke ogsaa deres Lys for at der kan
herske det forønskede Mørke. Man tage hvilke Begreber man vil: Princip, Lov,
Consequens --- de anvendes saaledes, at de Intet skulle udtrykke af deres
Væsens Indhold, ja at de \emph{Intet} skulle udtrykke, for at de kunne betegne
det, deres Tyranner ere; thi disse ere Intet uden dem, de have underkuet, og
blive, naar de have underkuet dem, dog ikke Andet end Jammerlighed: af Andre
maae de laane Alt, --- som man siger, indtil Skjorten --- og saa give de sig
Mine af at være det aandelige Livs
% NO FULL STOP after `aandelige'.  The Google layer reads `det aandelige. Livs'
% and KB reads it solid; the mark was measured at 900 dpi.  It is 12x14 px,
% area 115, and its foot sits 10 px ABOVE the baseline; the certain full stop
% after `udtrykke' in line 1 of this page measures 22x22, area 307, with its
% foot ON the baseline.  It is a speck of foul ink, not a period — the same
% class of Google error as the invented stop at p. 59 (ORTHOGRAPHY.md §3e).
% `det aandelige Livs egentlige Capitalister' is also the only reading that
% construes.
egentlige Capitalister, eller rettere sagt, saa udgives de derfor.

I det Øieblik man virkelig tager Læren om det Antirationelle paa Ordet, saa
standser al Discussion, ja al Meddelelse; det skeer ikke, fordi man ikke rigtig
har forstaaet Prof.\ Nielsen, men netop naar man rigtig har forstaaet ham: man
kunde ikke misforstaae ham stærkere end ved at give sig til at drøfte Sagen med
ham, han kunde ikke sætte en Begynder paa en vanskeligere Prøve end ved at
opfordre ham til en fælles Drøftelse. At det forholder sig saa, skal jeg
forsøge paa at vise, naturligvis vel vog\-%
% `naturligvis' — single i, as ORTHOGRAPHY.md §1 records for p. 48.
% --- p. 49 ---
% The word `vog-/tende' is broken across the page turn; resolved with \-%.
% Printer's signature `4' stands alone at the foot of this page: it is not text
% and is not transcribed.
% No footnote.  No letterspacing: spacing.py detected none and none is visible;
% KB's `Frim urer', `E r', `Han del', `Frim u reriet', `Ud tryk' are all ABBYY
% splitting at line ends and at ordinary justification gaps.
\opage{49}tende mig for at falde i den Snare, Sagen ifølge sin Natur lægger for
Enhver, der indlader sig med den.

For det Første er enhver Meddelelse umuliggjort: Naar en Frimurer vil prøve, om
en Anden er Frimurer, saa spørger han ikke den Anden derom, siger ikke heller,
at han selv er det, han bruger overhovedet ikke Ord for at give og modtage
denne Meddelelse, han kan tale om Vind og Veir, om Handel og Politik, ja han
kan endogsaa tale om Frimureriet, men vil han opnaae sit Maal, bruger han et
Meddelelsesmiddel, der ligger ganske udenfor Sproget: han seer at faae
Leilighed til at give den Andens Haand et eiendommeligt Tryk. Er den
Paagjældende nu ikke Frimurer, saa mærker han blot et Tryk eller han mærker
Intet: er han Frimurer, saa forstaaer han Meddelelsen og meddeler sig igjen.
Paa lignende Maade forholder det sig med den antirationel-religiøse Ethiker;
han kan lige saa lidt som Frimureren meddele sig igjennem Sproget, der jo er et
System af Tegn for Tanke- og Fornuftbestemmelser, men da han paa den anden Side
ikke som Frimureren kan bruge Haanden istedenfor Ordet, vælger han Sprogets Ord
til Betegnelse for Noget, der er ligesaa uafhængigt af de Fornuftbestemmelser,
de betegne, som Frimurerens Meddelelse er uafhængig af det Haandtryk, han
giver. Dette Noget er jo nemlig kun til for Organet for det
Antirationel-Religiøse, men er slet ikke til for den Fornuft, for hvilken
Sproget er Udtryk og Meddelelsesmiddel. Kan hint Noget nu ikke
% --- p. 50 ---
% ONE footnote, mark *), set at the closing quote of `Side“'.  Read for
% letterspacing BY EYE at the image (spacing.py is unreliable in footnote type,
% BATCH-AGENT.md §5): the note is set solid throughout.  It does NOT run over
% onto p. 51, which carries a footnote of its own.
% Two em-dashes round `thi ogsaa det er en Hemmelighed', taken from the image.
% The quotation is „…“ low-high, tight against `Side' with no space before the
% footnote mark; the Google layer sets it `,dialektisk … Side " *)'.
% No letterspacing in the body either.
\opage{50}meddeles gjennem Sproget, saa kan det naturligvis ligesaa lidt
meddeles ved noget andet Middel, der staaer under det Rationelles Lov;
Meddelelsens Vei er en Hemmelighed for Enhver, der ikke besidder Organet. Naar
man spørger, hvorfor da Sproget bruges, da maa dette uden Tvivl --- thi ogsaa
det er en Hemmelighed --- betragtes som et Incitament for at den hemmelige
Meddelelse kan træde i Kraft; spørger man derimod, hvorfor netop Sproget
bruges, da er dette uforstandigt spurgt, eftersom det samme Spørgsmaal kunde
gjøres med Hensyn til ethvert andet rationelt Middel og om de hemmelige Midler
kan der ikke spørges, da Spørgning er en rationel Act.

Lad os see Rigtigheden heraf i nogle Exempler. Naar der f.\ Ex.\ er Tale om en
antirationel-religiøs Ethik, saa betegner Ordet Ethik i den
Antirationel-Religiøses Mund og for den Antirationel-Religiøses Øre vel den
Fornuftbestemmelse, som Ordet udtrykker, men for det antirationel-religiøse
Organ betegner det noget ganske Andet, som vi, netop fordi det falder udenfor
al Fornuft, her umuligt kunne betegne. Naar det videre siges, at den paradoxe
Ethik „dialektisk forstaaet har sin rationelle
Side“\fnmark{*)}{Heegaard anf.\ Skr.\ p.\ 439.}, saa kan for det Første Ordet
dialektisk ikke betegne den Fornuftbestemmelse, for hvilken det er Udtryk, men
det betegner en Eiendommelighed ved det Antirationelle, der kun kan fattes af
Organet og som dette
% --- p. 51 ---
% ONE footnote, mark *), set immediately after `hedder,'.  Read for
% letterspacing BY EYE: set solid.  It does not run over onto p. 52.
% Printer's signature `4*' stands alone at the foot, below the note: not text.
% `Incitament.' — FULL STOP before `Naar', read at the image; KB reads a comma
% (`Incitament, N aar') and is wrong.
% No letterspacing on the page.  KB's `E t og Alt', `Betyd ning', `N aar' are
% ABBYY artefacts, rejected at the image.
% `naturligvis' — single i, as ORTHOGRAPHY.md §1 records for p. 51.
\opage{51}optager i sig, foranlediget derved, at Fornuften opfatter hvad Ordet
udtrykker, for det Andet kan Ordet rationel paa ingen Maade have sin vedtagne
Betydning, men det antyder kun en ny Eiendommelighed, for hvis Optagelse af
Organet Bestemmelsen det Rationelle er Incitament. Naar det endelig
hedder,\fnmark{*)}{Heegaard p.\ 406, 439, 445.} at det Antirationelle i Ethiken
fremkommer som en Consequens af det Rationelle, saa betegner Ordet Consequens
her naturligvis ikke, hvad det synes at betegne, men endnu en ny og interessant
Eiendommelighed ved det Antirationelle, til hvis Optagelse
Fornuftbestemmelsen Consequens inciterer. Ja selv om Udtrykket det
Antirationelle gjælder dette. Kun for Tanken og Fornuften er det antirationelt;
det, for Opfattelsen af hvilket det er Incitament, kan altsaa ikke være det
Antirationelle, men er Noget, om hvis Forhold til, om hvis Gyldighed eller
Ugyldighed for Tanken der ikke kan være Tale, da dets Eiendommelighed i Et og
Alt kun er for Organet. Og saaledes overalt: Ordene med deres
Fornuftbestemmelser ere ligesaa mange Incitamenter, der foranledige det
antirationel-religiøse Organ til dets mange særegne Functioner, for hvilke
disse Ord aldeles ikke ere Udtryk; ihvad hine end ere, meddeles kunne de ikke.

Det Samme, der er Grunden til, at det Antirationelle ikke lader sig meddele,
begrunder ogsaa, at det ikke kan objectiveres og overhovedet ikke tænkes
% --- p. 52 ---
% No footnote, no letterspacing (spacing.py: none; confirmed at the image).
% `„absolut“' — low-high quotes, set TIGHT, with the comma outside the closing
% mark.  Google reads `„absolut “' with a space; that is the justification
% artefact of ORTHOGRAPHY.md §3f and KB reads `„absolut"'.  Both are wrong.
% Two em-dashes round `der forholder sig til det Ideelle', from the image.
\opage{52}eller gjøres til Gjenstand for nogen Erkjendelse, i hvilken
Tænkningen er et Moment. Grunden er den, at det falder absolut udenfor al
Rationalitet. At der maa lægges Eftertryk paa „absolut“, er en Selvfølge, da
Intet i Tilværelsen gaaer op i Rationalitet, da det Rationelle overalt
forudsætter det Reelle, uden hvilket det selv ingen Realitet besidder. Men
Rationelt og Reelt gjennemtrænge hinanden, saa de høre til hinandens Væsen,
eller bestemtere: det er samme Væsen, der adskiller sig i sine to Væsenssider.
Det Rationelle staaer da under det Reelles Lov og det Reelle under det
Rationelles. Viser der sig derfor et Reelt, der ganske unddrager sig det
Rationelles Lov, saa kan selvfølgelig ingen rationel Bestemmelse i mindste
Maade udtrykke dets Væsen, og det kan hverken være til for Tænkningen --- der
forholder sig til det Ideelle --- eller for den Erkjendelse, der svarer til det
Reelle, fordi Tænkningen overalt er et Moment i den, ligesom det Ideelle
overalt er et Moment i det Reelle. Bestemmelsen det Reelle kan da lige saa lidt
som nogen anden anvendes paa det Antirationelle, den er kun som hine tidligere
nævnte Incitament for, at Organet kan fungere i en vis bestemt Retning. Det
Samme gjælder saa om Ordet fungere og de øvrige Udtryk, der alle betegne Noget,
hvori er et Rationelt.

Hvis det Antirationelle ikke opfattes saaledes, saa forvandles det til et
Rationelt; thi et Rationelt er et Noget, der er reelt, et Reelt er et Noget,
der er
% --- p. 53 ---
% ONE letterspaced word: `kjende' in `Vi k j e n d e altsaa ikke noget
% Antirationelt'.  Read at the image: `Vi' immediately before it is SOLID, and
% so is the whole of the following line (`altsaa ikke noget Antirationelt,
% kjende intet Organ'), which was checked because the second `kjende' five
% words later is the same word unemphasised.  spacing.py flagged `kjende' only,
% and the image agrees.
% KB's `h ø re r' (`hvortil da hører') and `e r' (`tirationelle er, det er')
% were both checked at the image and are SOLID — ABBYY noise, not Sperrsatz.
% One em-dash after `aldeles ikke udtrykke', from the image.
\opage{53}rationelt, men et Antirationelt er et Noget, der hverken er reelt
eller rationelt, der ikke er et Antirationelt, der vel kan kaldes saa, men dog
egentlig ikke kaldes saa med Rette, da det aldeles Intet kan have at gjøre med
det Udtryk, hvormed det benævnes. Deraf følger saa tillige, at der ingen
Videnskab kan dannes om det; der kan meget vel dannes et System af Tanker, men
dette vil blot sige et System af Incitamenter for det Antirationelles
hemmelighedsfulde Virken, som hine Tanker aldeles ikke udtrykke --- hvortil da
hører, at Ordet Incitament heller Intet udtrykker, der angaaer Forholdet mellem
det Rationelle og det Antirationelle, da der slet intet Forhold er. At der
ingen Kritik kan gives af det Antirationelles Fremstilling, da der ingen
Fremstilling kan findes af det; at ingen Discussion kan tænkes med den, der
staaer paa det Antirationelles Standpunkt, behøver man nu kun at minde om,
eftersom Kritik og Discussion ere Acter, der staae under det Rationelles Lov.

Naar vi sluttelig erindre, at det at kjende er en rationel Act, hvilket ikke er
det Samme som en blot rationel, saa følger det af sig selv, at vi ikke kjende
noget Reelt, der ikke tillige er rationelt. Vi \emph{kjende} altsaa ikke noget
Antirationelt, kjende intet Organ for det Antirationelle. Er det da ikke? Nei;
thi det at være har en rationel Side; at sige, at det Antirationelle er, det
er: at sige, at det er rationelt.
% --- p. 54 ---
% Last page of chapter 3.  FIVE lines of text, continuing the paragraph begun
% on p. 53 (the first line is flush to the margin, not indented), then a short
% centred double rule closing the chapter — the same ornament that closes
% chapter 2 on p. 46.  The rest of the leaf is blank.  The page drops nothing
% else: no footnote, no signature, no letterspacing (spacing.py: none;
% confirmed at the image).
% p. 55 opens chapter 4 and is not part of this batch.
\opage{54}Hvad skulle vi da sige om det? Vi kunne Intet sige, ligesaa lidt
bekræfte som benægte Noget derom, og man maa fremfor Alt erindre, at hvad her
er sagt om det, Intet udtrykker om dets Væsen, er absolut og ubetinget fremmed
for det.

\begin{center}
\rule{2cm}{0.4pt}\\[-4pt]
\rule{2cm}{0.4pt}
\end{center}

%% ============================================================
%%  4. DUALISMENS CONSEQUENSER  (printed pp. 55--76 = PDF 59--80)
%%  Head as printed: `4.' / `Dualismens Consequenser.'
%%  Indhold reads `Dualismens Consequenser', without number and without stop.
%%  p. 76 closes the chapter and then carries an unheaded Efterskrift.
%% ============================================================
\chapter*{4.\ Dualismens Consequenser.}
\addcontentsline{toc}{chapter}{4. Dualismens Consequenser}
\markboth{Dualismens Consequenser}{Dualismens Consequenser}
\label{ch:4}

% --- p. 55 ---
% Chapter opening: `4.' / `Dualismens Consequenser.' is already set by the
% skeleton's \chapter*, so it is NOT repeated here.  The page drops its folio
% and opens with an enlarged capital `N' (two lines deep); per the batch brief
% the enlarged capital is not marked up.  spacing.py's `single?: Dualismens
% (1.24)' is the chapter head, which is set in a larger display face, not
% letterspaced --- read at the image.  No letterspacing anywhere on p. 55.
% TWO footnotes, marks `*)' and `**)'.  `Propædeutik' has the æ ligature: read
% at 600 dpi.  The Google layer reads `Propedeutik' (ORTHOGRAPHY.md §3b); KB
% reads `P hil. P ropæ deutik'.  Neither was copied.
\opage{55}Naar der er kommen en lille, næsten usynlig Revne i et Glas, saa kan
den undertiden, efter længe at have holdt sig paa samme Punkt, med Eet slaae
igjennem og spalte hele Glasset. Naar en af det menneskelige Legemes
Extremiteter angribes af en Sygdom, da kan Skaden fra dette Udgangspunkt æde
sig ind i Legemet og angribe dets Sundhed i sin Rod. Saaledes er det gaaet med
den Dualisme, der tidligt\fnmark{*)}{Phil.\ Propædeutik 1857 p.\ 77.} ytrede
sig i Prof.\ Nielsens Anskuelse af Forholdet mellem Tro og Viden. Den syntes
længe uskadelig for hans Theorie; thi den angik jo ikke Theorien selv, den
beroede jo kun paa, at denne Theorie ingen Indflydelse fik paa hans Praxis; men
nu har den slaaet saa kraftigt igjennem, at den er trængt ind i selve hans
philosophiske Grundanskuelse\fnmark{**)}{Bemærkninger om theoretisk og praktisk
Erkjendelse. Nord.\ Universitetstidsskr.\ Tillægshefte til 10de Aargang, 1866.
Kjbh.}. I en vis Forstand kunde dette synes forunderligt: man skulde troe, at
det ingen Virkning vilde have paa
% --- p. 56 ---
% One footnote, mark `*)'.  Its page reference is 343: read at 600 dpi.  The KB
% layer reads `p. 543' --- KB is wrong, as ORTHOGRAPHY.md §3g-bis predicts for
% footnotes.
% `helt ind i Videnskaben' --- SINGLE e, as printed (cf. `heelt' p. 61).
% The French is set in the ordinary roman, no contrasting face, so no markup:
% `comme une base, qui peut porter la trinité chrétienne'.  Both accents are
% acute and were read at the image.
% No letterspacing on this page (spacing.py: none; confirmed at the image).
\opage{56}Philosophens Lærebygning, at han udenfor sin Videnskab havde en
uvidenskabelig Forvisning. Selv om Prof.\ Nielsen vilde gaae saa vidt at
forbyde os at anvende den Theorie, han gav os, nuvel --- kunde man sige --- saa
lad os idetmindste være theoretiske i Theorien. Var Meningen ikke den, hvorfor
har da Professoren vel ellers saa tidt og med saa stor Fornøielse fortalt os
Anecdoten om Laplaces Svar paa Napoleons Spørgsmaal? Vilde det ikke være høist
forunderligt, om den Philosoph, der saa stærkt har betonet Videnskabens Protest
mod ethvertsomhelst Dogme, skulde ende med som Victor Cousin i Frankrig at
frembyde sit Gudsbegreb „comme une base, qui peut porter la trinité
chrétienne“?\fnmark{*)}{se Grundid.\ s.\ Logik II.\ p.\ 343.} Men Sagen er den,
at netop fordi den Adskillelse, som Prof.\ Nielsen vil anbringe mellem Livet og
Videnskaben er saa chimærisk, netop derfor strækker Existensens Spaltning, den
„psychologiske“ Splid sig helt ind i Videnskaben, hvorhen den aldeles ikke
skulde kunne naae. I sine „Bemærkninger om theoretisk og praktisk Erkjendelse“
har Professoren ladet Spaltningen fra det religiøse Omraade slaae over paa det
ethiske, og fra det ethiske paa Erkjendelsens. Her er det da ikke blot Religion
og Videnskab, men Handlen og Erkjenden og indenfor Erkjendelsen igjen den
theoretiske og den „praktiske“ Erkjendelse, der staae som absolut uensartede
overfor hinanden. Det har hermed
% --- p. 57 ---
% THE MARK `*)' IS PRINTED TWICE ON THIS PAGE AND THERE IS ONLY ONE NOTE at the
% foot (`Heegaards Fremstilling...'), both quotations being from Heegaard.  The
% note is set here at the first mark; the second occurrence, after
% `Genesis“', is transcribed as the literal `*)' the page prints, since it
% refers back to the same note.  Read at 600 dpi: the second mark is a SINGLE
% asterisk, not `**)'.  The footnote itself carries no letterspacing (checked
% by eye at 600 dpi; KB's `F rem stillin g', `L æ re' are ABBYY noise).
% FOUR letterspaced runs, all read at 600 dpi:
%   `Kraft af den samme Udviklingslov som gjælder indenfor det Rationelle'
%   `Fornuft'          (`bestemmer det Logiske' is NOT spaced)
%   `ved det Antilogiske'   (`og omvendt“' is NOT spaced --- checked)
%   `(et Antilogisk)'  --- the parentheses share the widened spacing
%   `Paradox'          --- in `„et rationelt Paradox“' ONLY the last word is
%                          letterspaced; `et ratio-/nelt' is set solid.
\opage{57}haft samme Forløb som med Paradoxet. Om et Paradox, et absolut
Paradox i og for sig kan man mene hvad man vil; men hvad man saa mener, maa det
dog vel være klart, at et formidlet Paradox er en fuldkommen Uting. Man vil
ogsaa erindre, hvor latterlig Kierkegaard fandt al „Skimten“ ligeoverfor det
absolute Paradox (Afsl.\ Eftskr.\ p.\ 433); men hvad er det vel Andet end den
aller morsomste „Skimten“ naar det hedder\fnmark{*)}{Heegaards Fremstilling af
Prof.\ Nielsens Lære, Indledn.\ til Ethiken p.\ 439.} at „det paradoxe
Standpunkt fremkommer i \emph{Kraft af den samme Udviklingslov som gjælder
indenfor det Rationelle}“ og at „den rationelle Videnskab fordobler sit
Synspunkt saaledes, at den paa eengang erkjender det Paradoxes existentielle
Ubegribelighed og just i Kraft af denne Erkjendelse tillige opdager den
Udviklingslov, hvorefter Videnskaben giver en Afspeiling af det Paradoxes
Genesis“*).
% `*)' above: the page's second, repeated reference to the note already set.
Saaledes kommer det relative Paradox til at forberede det absolute og bane
Veien for det. Medens det dog endnu i Grundid.\ s.\ Logik II.\ p.\ 278 hedder
at „Magten i sit Væsen er lige saa rationel som Viden,“ og at „den uendelige
\emph{Fornuft} bestemmer det Logiske \emph{ved det Antilogiske} og omvendt“
bliver hos Dr.\ Heegaard selve det at ville, forsaavidt det er et for den
theoretiske Viden Incommensurabelt (\emph{et Antilogisk}) kaldt „et rationelt
\emph{Paradox}“ (p.\ 441) og dette vakre Begreb
% --- p. 58 ---
% TWO footnotes, marks `*)' and `**)'; the second reads `P. 66.' with a CAPITAL
% P, read at 600 dpi.  (The batch brief warned the Google layer might be silent
% about the second note; in this extraction it is not --- it reads `** ) P. 66.'
% --- and the image carries both.)  Neither note is letterspaced.
% ONE letterspaced run: `ikke ligefrem' in `at Praxis ikke ligefrem lader sig
% aflede af Theorie'.  The SECOND `ikke ligefrem' on this page (`at
% Videnskaben ikke ligefrem kan modbevise Miraklet') is set solid --- read at
% the image, and spacing.py flags only the first.
% `helt ned i Immanensens' --- single e, as printed.
% `Prof.:' --- the colon is on the image and in KB; the Google layer drops it.
\opage{58}danner saaledes den forønskede Overgang. Som det Paradoxe nu her har
skudt Rødder helt ned i Immanensens rationelle Verden, saaledes ogsaa den
absolute Uensartethed.

Det er imidlertid vanskeligt nok paa dette Punkt at faae en tilstrækkelig skarp
og afgjort Udtalelse af Prof.\ Nielsen. Han bruger vel stærke Ord, men han
tilføier gjerne en lille Bisætning eller Bibestemmelse, der atter berøver
Ordene deres „resolute“ Præg, han tager med den ene Haand, hvad han har givet
med den anden. Han befæster Gabet mellem Praxis og Theorie, han synes endog at
gaae saa vidt, at han vil negte, at overhovedet Theorien som saadan lader sig
overføre i Praxis; thi --- hedder det --- „det, der erkjendes theoretisk, kan
ikke erkjendes praktisk“\fnmark{*)}{Anf.\ Sted p.\ 69.}. Denne Paastand synes
dog at være af en Uforsigtighed, der idetmindste har det for sig at være
„resolut“. Men rettede nu En det Spørgsmaal til Prof.: Siger Du da virkelig, at
Praxis ikke lader sig aflede af Theorie? da kan han svare med sine egne Ord:
Ingenlunde, jeg har kun sagt, at Praxis \emph{ikke ligefrem} lader sig aflede
af Theorie\fnmark{**)}{P.\ 66.}. Og hvo i al Verden vilde modsige det? Paa
samme Maade erklærede Hr.\ Professoren i Indledningen til Grundid.\ s.\ Logik
I, at Videnskaben ikke ligefrem kan modbevise Miraklet. En værdifuld Kategori,
den
% --- p. 59 ---
% One footnote, mark `*)'.  `Propædeutik' again with the æ ligature, read at
% 600 dpi; the Google layer reads `Propedeutik'.  Unlike the p. 55 note, this
% one prints a full stop after `Propædeutik'.  No letterspacing in the note
% (read by eye at 600 dpi, since spacing.py is unreliable in footnote type).
% ORTHOGRAPHY, all three read at the image and NOT normalised in either
% direction: `deels ... deels' (p. 42 prints `dels'), and `naturligviis' ---
% the ONLY one of the book's twelve occurrences that doubles the i.  The Google
% layer silently modernises it (ORTHOGRAPHY.md §1); KB has it right.
% `Fremdeles:' HAS a colon --- on the image and in KB; Google drops it.
% `Oprindelighed med Tanken (p. 68)' --- no stop after `med'; Google's `med.'
% is a speck promoted to a full stop (ORTHOGRAPHY.md §3e).
% ONE letterspaced run: `ikke ligefrem' in line 1.  `den absolute Uensartethed
% var bevist' is NOT letterspaced --- the wide gaps there are justification,
% the letters within each word are tight (read at 600 dpi).
\opage{59}Kategori: \emph{ikke ligefrem}, værdig den Mand, der saa skarpt har
bebreidet den speculative Dogmatik „den svampagtige Blødhed, hvormed Kategorien
Grad overalt i ethvert nok saa afgjørende Dilemma, deels direkte, deels
indirekte maa figurere“\fnmark{*)}{Phil.\ Propædeutik. Efteraarshalvaar 1865
p.\ 173.}.

Fremdeles: Prof.\ sprænger Broen i Luften mellem Tanken og Villien, efter
først at have slettet Følelsens Territorium ud af sit Landkort og delt denne
svagere Magts Besiddelse mellem hine eneraadige Magthavere. Han tillægger ikke
blot Villien lige Oprindelighed med Tanken (p.\ 68), men lægger en saadan Vægt
paa Gjendrivelsen af den Hegelske Vildfarelse, at Villien kun er en
Modification af Tænkningen, at man skulde troe, den absolute Uensartethed var
bevist, saasnart blot denne Misforstaaelse var gjendrevet. Men Spørgsmaalet er
jo her slet ikke det, om Villien lader sig aflede af Tanken eller omvendt, men
om begge disse Væsensytringer ikke lade sig aflede af et og samme Grundvæsen.
Negtes dette, da have vi Dualismen, ren og skjær. Men det Afgjørende eller
rettere Mangelen paa Afgjørelse er ogsaa her, at det i Virkeligheden kun
tilsyneladende benegtes. Thi indvender man, at Tanken og Villien dog staae i
uafbrudt Forhold til hinanden, da svarer Prof.\ blot: det har jeg jo aldrig
lagt Skjul paa „men Tanken er derfor ikke Villie, Villien ikke Tanke o.\ s.\
v.“ (p.\ 67) hvad det naturligviis ikke kunde falde noget
% --- p. 60 ---
% No footnote on this page.
% ONE letterspaced run: `grændsende', broken across the line as `g r æ n d-' /
% `s e n d e' and read at the image.  spacing.py flags both halves separately.
% ORTHOGRAPHY: `anden Del p. 457' --- single e, as printed (cf. `Deel' p. 70),
% and `naturligvis' single i in the quoted passage (cf. p. 59 above).
% The quotation `„En mat Idee ...' OPENS here and does not close until
% `Slutningen.“' on p. 61: the „ on this page is deliberately unmatched.
\opage{60}Menneske ind at benegte. Forsøger man da for at naae tilbage til
Centrum at vende Prof.s egne Ord imod ham, som følger: Hvis „Religionen
forudsætter en Aabenbaring og Aabenbaringen et Under, men Videnskaben fornegter
Underet og forkaster Aabenbaringen“ (p.\ 63) saa kan den Form af det Religiøse,
der lader sig forene med Videnskaben, kun være en saadan, der er uden Undere og
uden Aabenbaring --- da svarer Prof.\ Nielsen bittert, at „slige
Expectorationer mere charakterisere de livlige end de grundige Hoveder“ (p.\
73) uanseet, at det er ham selv, der har opstillet Præmisserne, af hvilke denne
Slutning med Nødvendighed maa drages. Opgiver man da Modet og indvender blot
tilsidst, at Prof.s Anskuelse da maa være en aabenbar Dualisme, kan
„Bemærkningernes“ Forf.\ med Rette svare: Jeg siger det jo selv, jeg er
næstendels Dualist; thi --- hedder det --- „der gives i Aandslivet stærke til
Dualisme \emph{grændsende} Modsætninger“ (p.\ 66). Næstendels-Dualismen er
beslægtet med den af Prof.\ N.\ saakaldte „Middelveis-Theologie“; man fanger
den ikke, den er saa smidig at den glider Enhver, der vil gribe den, afhænde.

Det er skjønne, sande og begeistrede Ord, som Grundideernes Logiks Forf.\
udtaler i anden Del p.\ 457: „En mat Idee har naturligvis matte Extremer med
svag Midte og mat Sammenslutning, derfor er den immanente Speculation altid saa
ængstelig for hvad der kunde ligne en Dualisme. Den mægtige, i
% --- p. 61 ---
% One footnote, mark `*)' (the asterisk is ink-damaged on this copy but is a
% single star at 600 dpi).  No letterspacing in the note.
% ORTHOGRAPHY: `heelt ud' --- DOUBLE e, as printed, the book's one `heelt'
% against `helt' on pp. 27, 40, 56, 58.  Read at the image.  Do not normalise.
% No letterspacing in the body of this page (spacing.py: none; confirmed).
% The closing “ in line 4 closes the quotation opened on p. 60.
% KB reads `relative”' with the wrong closing mark; the image has „…“.
\opage{61}Væsen guddommelige Idee har derimod mægtige Extremer, dens Udslag i
kraftig Dualisme er kun Forudsætning for den uendelige Overlegenhed, hvormed
den fuldfører Slutningen.“ Men hvad hjælper det os, at disse Ord ere sande,
naar i det Psychologiske „Slutningen“ udebliver og uforsonlige Modsigelser
stilles Ansigt til Ansigt mod hinanden. Vi skylde alle Prof.\ Nielsen Tak og
Anerkjendelse for den Energi, hvormed han har forfulgt det Ideelle heelt ud i
dets Modsætnings yderste Extrem, Beundring for den geniale Tillid, hvormed han
altid har følt sig forvisset om, at det vel skulde lykkes at vise selv det
Stridigstes Eenhed i den mægtige Idee; vi høre ham uden Uvillie tale haant om
dem, der vilde skræmmes selv af Skyggen af en Dualisme; men, det forstaaer sig,
dog kun saa længe, som den virkelige Dualisme ikke bliver udgivet --- som
Manden i Eventyret --- for sin egen Skygge. Er Tilliden hos Læreren først
bleven til Overmod, da svinder hos Lærlingen Trygheden for Følelsen af en
virkelig Søndersplittelse, en Spaltning, der ikke lader sig skjule under
ubestemte Vendinger, som de nylig anførte, et Brud, der ingenlunde lader sig
klinke ved den Art Dialektik, der bestaaer i en Sammentvingen af modsigende
Udtryk som „rationelt Paradox“ eller „absolut uensartede Ytringer, der vise sig
relative“\fnmark{*)}{Bemærkningerne p.\ 65.}.

Heldigvis lader „Bemærkningernes“ Forf.\ sig ikke
% --- p. 62 ---
% No footnote on this page, and no letterspacing: spacing.py's `single?: Sag'
% and `single?: der' were both rejected at the image (ordinary type, wide
% justification gaps only).
% `væsentlig praktisk' here has the t; p. 63 prints `væsenlig'.  Both stand as
% printed.
% The quotation `„Vistnok falde ...' opens here and closes at `Bevis.“' on
% p. 63, so the „ on this page is deliberately unmatched.
\opage{62}nøie med at udtale almindelige Sætninger, han vælger bestemte
Exempler til Oplysning af det Fremsatte; men som det i en saa mislig Sag var at
vente, Prof.\ N.\ er sært uheldig med de Exempler, han giver. Han udtrykker sig
(p.\ 69) paa følgende Maade: „Vistnok falde saadanne Kategorier som Valgfrihed,
Valg, Beslutning o.\ fl.\ ind under Tankens Belysning, forsaavidt de klare sig
i sædelig Erkjendelse; men i Principet ere de ligefuldt Villies-Kategorier,
deres Erkjendelse er væsentlig praktisk. Det, der erkjendes theoretisk, kan
ikke erkjendes praktisk; det nytter ikke at beraabe sig paa, at Indholdet, der
erkjendes, er det Samme, kun Erkjendelsesmaaden forskjellig; i enhver
afgjørende Erkjendelse bliver Indholdet bestemt ved Maaden, og Maaden ved
Indholdet. At Modet som Dyd kan erkjendes i Theorien, maa indrømmes; --- men
dersom det Mod, der erkjendes theoretisk, virkelig i Eet og Alt var det samme,
det selvsamme, som det, der erkjendes praktisk, hvad vilde saa heraf følge?
Heraf vilde jo følge, at en dygtig theoretisk Tænker, uden at vove Noget, uden
at udsætte sit Liv for nogen virkelig Fare, maatte kunne tænke sig saaledes ind
i det sande Mod, at han i Virkeligheden med det Samme tænkte det sande Mod ind
i sig, og saaledes ved objectiv Tænkning, hvor feig han for Resten var af
Naturen, virkelig blev modig. Men at den, der i Virkeligheden er feig, om han
end tænker nok saa theoretisk, ligesaa lidt kan tilspeculere sig virkeligt Mod,
som den, der er uden Penge, om han
% --- p. 63 ---
% No footnote, no letterspacing (spacing.py: none; confirmed at the image).
% `væsenlig er praktisk' --- WITHOUT the t, as printed here, against
% `væsentlig' on p. 62.  Read at the image; both layers agree.  Not normalised.
% The closing “ in line 3 closes the quotation opened on p. 62.
% `Bemærkningerne“,' --- KB reads `gerne11' (its usual garbling of “).
\opage{63}end vilde tænke theoretisk objectivt, til han revnede, kan
tilspeculere sig virkelig Penge, trænger ikke til Bevis.“

Jeg skal ikke videre opholde mig ved det sidste Exempel med Pengene, som man
ikke skal kunne tilspeculere sig, om man saa tænkte, indtil man --- efter
Prof.\ N.s kraftige Udtryk --- revnede; jeg troer, man uden Chicane kan
paastaae, at enhver Kjøbmand, ja endog blot enhver Børsspiller, der har
speculeret, vil kunne lære Prof.\ Nielsen, at man kan speculere sig Penge til,
og Penge fra. Eller er det maaskee i Ordene „theoretisk objectivt“ at Pointen
stikker? Jeg troer dog knap, at Børsspeculationer kunne henregnes til den
subjective, den ethisk-praktiske Tænkning. Men hvad jeg vil holde mig til, er
Exemplet med Modet. Thi dette Exempel er et Vidnesbyrd fremfor alle om den
Mangel paa Klarhed og besindig Eftertanke, der charakteriserer hele den nye
Lære. Først gjøres der en Inddeling af Kategorierne i Villieskategorier og
Videnskategorier; saa siges der om de første, at deres Erkjendelse væsenlig er
praktisk; dernæst nævnes Modet som Exempel. At Prof.\ Nielsen lægger Vægt
herpaa, det sees deraf, at dette Exempel ikke alene forekommer een Gang i
„Bemærkningerne“, men hos Echoet Hr.\ Dr.\ Heegaard fulde fire Gange, nemlig
p.\ 385, 399, 400, 405. Hvis her nu med Modet mentes den Uforsagthed, der
næsten som et physisk Anlæg kan være eet Menneske medfødt og et andet berøvet,
da vilde Den visselig med Rette
% --- p. 64 ---
% No footnote on this page.
% CAPITAL Æ: `Æren kræver' --- the ligature read at 600 dpi (ORTHOGRAPHY.md
% §3a lists p. 64 among the six places the Google layer loses it; it reads
% `Eren' here, KB reads `Æ ren').  Every other word-initial capital vowel in
% pp. 55--65 was checked at the image as well (Eet, En, Er, Erkjendelse,
% Exempel, Extremer, Et, Alt, Andet, Aandens, Anlæg) --- all plain E/A.
% PRINTER'S ERROR, transcribed as printed: `Den thoretiske Erkjendelse' for
% `theoretiske'.  Read at 600 dpi; both layers agree on `thoretiske'.
% `tungnem' (not Google's `tunguem') and `overvinder' (not KB's `oven inder')
% read at the image.
% THREE letterspaced passages, all read at the image:
%   `praktisk'  in `Modet som praktisk Bestemmelse'
%   `Dyd'       in `Modet som Dyd'
%   `hvor ivrigt han saa bestræber sig for at gjøre det til Et og Alt.'
\opage{64}være latterliggjort, som troede, at han, hvor feig han end af Naturen
var, kunde ved objectiv Tænkning blive modig --- af Naturen. Men hvad da blev
sagt om Modet som \emph{praktisk} Bestemmelse, vilde rigtignok uheldigvis ogsaa
gjælde i samme Maal om alle de theoretiske Anlæg; thi man kan ligesaa lidt,
naar man af Naturen er tungnem, sentænkende og svagtbegavet tilspeculere sig
medfødt Lærenemhed, Skarpsindighed eller Vid. Men da her nu, som det siges,
udtrykkeligt sigtes til Modet som \emph{Dyd}, er den Bestemmelse som gives,
ligefrem falsk: Den thoretiske Erkjendelse af, hvad Sandheden byder, hvad et
Pligtforhold er, hvad Æren kræver, hvad et Fædreland vil sige, kan og skal
gaae saaledes over i Kjød og Blod, at den bekjæmper og overvinder den naturlige
Forsagthed, hvor denne staaer Udførelsen af det Rette imod. Hvo der negter
Muligheden heraf, han burde ikke driste sig til at føre Livets Sag contra
Theorierne; thi han røber sig selv som Stuephilosoph, som ren Theoretiker,
hvormeget han saa lovpriser, ja forguder det Praktiske, \emph{hvor ivrigt han
saa bestræber sig for at gjøre det til Et og Alt.}

Thi ret beseet, hvad gjør Prof.\ Nielsen Andet? Han taler vel om den dobbelte
Verden, han synes vel tidt at være ren Dualist, men i sidste Instans vil man
finde, at det gaaer ham som alle Dualister, de ere det nødvendigvis kun
tilsyneladende, de nødes, naar det kommer til Stykket, til at bortkaste det ene
af de tvende modstaaende Led, som de have bragt i
% --- p. 65 ---
% No footnote.  The numeral `5' standing alone at the foot is printer's
% signature 5 (the signatures fall at pp. 17, 33, 49, 65), not text.
% TWO letterspaced passages, read at the image:
%   `Aandens Liv er i Villien'  (the closing “ and stop are not spaced)
%   `forklare'  in `vil forklare os Tilværelsen'
% `thi naar galt skal være kan hin Side opofres' --- as printed: the gap after
% `være' is a justification gap, there is NO dash and NO comma on the image
% (checked at 600 dpi).  A comma may have dropped out in the forme; not
% supplied.
% FRENCH, set in the ordinary roman with no contrasting face, so no markup.
% Every accent read at 600 dpi: `siècle' and `Où' and `mise à' grave;
% `volonté', `idées', `isolées', `faculté', `élevée' acute; `l'âme', `l'être'
% circumflex.  `francais' is printed WITHOUT the cedilla --- as printed.
% `XIXe' sets a raised e.  After `volonté' the page prints FOUR spaced dots;
% set here as \dots followed by a stop, which is the nearest LaTeX equivalent.
% PRINTER'S INCONSISTENCY, transcribed as printed: the long French quotation
% is CLOSED at `est moi“' but was never opened --- there is no „ before
% `Il n'observe' on the image.  This leaves the fragment with one unpaired
% closing “; it is the page's, not the transcriber's.
% Also as printed: `par analyse.' takes a FULL STOP (round, no tail at 600 dpi;
% KB reads a comma) and the next sentence opens with a lower-case `c'est'.
\opage{65}unaturlig Spænding mod hinanden. Det viser sig, at Prof.\ N.,
polemiserende mod det at gjøre Tanken til Alt og Tænkningen til Eet med det
hele Menneske, (p.\ 68) i Virkeligheden gjør en anden af Menneskeaandens
Ytringer, nemlig Villiens Liv til Alt, til det Hele; thi naar galt skal være
kan hin Side opofres, men ikke denne (p.\ 76). Man høre Forf.s egne Ord: „Den,
der sætter Tanken ind paa praktisk Erkjendelse af en hellig Villie, gaaer
maaskee glip af Theorien, men han vinder Livet, Aandens evige Liv; thi
\emph{Aandens Liv er i Villien}“. Er det en Philosoph eller er det en Præst, en
Tænker eller en Moralist, som taler? Af den der taler med en sædelig eller
gudelig Tendens for Øie, finder man et saadant Udtryk berettiget og billiger
det gjerne; men af den, der vil \emph{forklare} os Tilværelsen som Tænker?

Man burde maaskee undre sig mindre derover, siden man ogsaa i andre
Litteraturer træffer paa ganske tilsvarende Misvisninger. Se f.\ Ex.\ Taine:
„Les philosophes francais du XIX\textsuperscript{e} siècle“ om Maine de Biran
p.\ 59: Il n'observe que le moi, l'âme, l'être, la substance. Où les
trouve-t-il? Dans la volonté\dots. Ainsi le moi n'est plus ce tout indivisible
et continu, dont nos idées, nos plaisirs, nos peines sont les parties
composantes, isolées par fiction et par analyse. c'est une force ou faculté,
portion du tout, mise à part, élevée au-dessus de toutes les autres. Le reste
est mon bien: celle-ci est moi“. Paa denne Maade charakteriserer den aandrige
Franskmand denne Lære, der gjør
% --- p. 66 ---
% Continues from p. 65 at a clean word boundary: p. 66 opens `Villien til det
% hele Menneske.'  KB reads `Yillien' (its characteristic V->Y); the image has
% the V.
% Every word-initial capital vowel in pp. 66--76 was read at the image for the
% Google-loses-Æ class (ORTHOGRAPHY.md §3a).  Two genuine Æ found, both as the
% audit predicted: `Æren' p. 70 and `Ære' p. 73 (see the comments there).  All
% the rest are plain E/A: Eet, Er, En, Een, Endnu, Existensen, Erkjendelse,
% Eiendommelighed, Eenhed, Efterskrift, Element, Ethik, Andet, Aanden, Amphi-
% bier, Afkald, Afstand, Aandløshed.
% No letterspacing on this page (spacing.py: none; confirmed at the image).
% Both em-dashes on this page were taken from the image: Google renders the
% first as a hyphen and the second as `-—', KB gets both right.
\opage{66}Villien til det hele Menneske. Men skjøndt saaledes Phænomenet ikke
ere enestaaende, maa jeg dog bekjende, at jeg for min Del blev høilig
% `Del' here, NOT `Deel' --- p. 70 below prints `Deel'.  Both read at the
% image; the book is inconsistent with itself and is not normalised
% (ORTHOGRAPHY.md §2).
forbauset, da jeg i Slutningen af Prof.\ N.s Afhandling stødte paa det Sted,
der begynder saaledes: „Skulde det endnu virkelig hvad Nogle paastaae være den
ene og samme Bevidsthed umuligt at anerkjende baade det Praktiske og
Theoretiske, Tro og Viden,“ --- man vil erindre at Prof.\ med disse Ord
betegner den positive Religion og Menneskeaandens intellectuelle Fordring ---
„da vilde det i alt Fald blive det Theoretiske, ikke det Praktiske, der maatte
opgives“. Hvorledes? Sagen er da altsaa ikke sikret og bevist, det kan endnu
være muligt, at Dualismen ikke lader sig gjennemføre, at man maa opgive at
sige baade Ja og Nei? Og jeg, som troede at denne Opfattelse --- efter Forf.s
Mening --- var aldeles gjendrevet, som manglende de nødtørftigste
Forudsætninger i Dialektik og Psychologie (se p.\ 65). Denne pludselige
Indrømmen af det Modsattes Mulighed har noget meget Beroligende. Det gjælder
da nu om at forandre Taktik, det er ikke mere Dualismen der skal angribes, det
er Udveien ud af den, som det bliver Opgaven at spærre. Denne Udvei findes nu
antydet i forskjellige Udtryk hos Prof.\ Nielsen og hans Tilhængere,
Bevægelsen forekommer i forskjellige Former, men i Virkeligheden er den let at
gjenkjende som een og den samme.

Paa dette Sted hedder det først og fremmest, at skal en af Siderne opgives,
saa maa det blive den
% --- p. 67 ---
% No footnote.  No letterspacing: spacing.py flagged a `single? ,er' (1.23),
% which is the opening „ of „er praktisk' read as a comma --- ordinary spacing
% at the image.  The signature `5*' at the foot is not text.
\opage{67}theoretiske. Man mærker vel, at Skibet er i Havsnød, det kan kun
frelses ved at kaste sin Ballast overbord: hvis det derved kun ganske sikkert
opnaaer at kæntre, saa er dets Undergang gjort saa temmelig indlysende. I ren
Almindelighed betragtet har nu den valgte Udvei noget meget Komisk. Hele
Læren bestod jo væsenligt deri, at man tog Kierkegaards „Øieblikket“ i den ene
Haand, Ørsteds „Aanden i Naturen“ i den anden, klappede dem sammen og indbandt
dem i Eet. Optræder nu En med den Indvending, at denne Forbindelse eller
Sammenbindning ikke godt gaaer an, og er han paastaaelig nok til ikke at ville
give efter, nu vel! saa vil Modparten ikke være den stridige, men skjærer
Ørsted fra og beholder Kierkegaard. Men hvad bliver der saa tilbage af den
hele Lære, som man paa saa taabelig en Maade har betegnet som den Mands
største Fortjeneste, hvis velerhvervede Berømmelse beroer paa ganske andre og
anderledes ægte videnskabelige Bedrifter, hvad indeholder den da, som ikke
langt bedre, rigere, fuldere er sagt af Kierkegaard?

Underkaste vi nu Sagen en nærmere Betragtning, hvorledes beviser Prof.\
Nielsen da Muligheden af at opgive det Theoretiske? „Existensen,“ hedder det,
„er praktisk; en Bevidsthed kan existere uden at være theoretisk; men en
Bevidsthed kan ikke være theoretisk uden at existere. Det er altsaa muligt at
sætte Existensen øverst og beholde Theorien.“ Er denne Sophisme ogsaa passende
for en Tænker, en Tænker
% --- p. 68 ---
% No footnote, no letterspacing.
\opage{68}af Prof.\ Nielsens Rang? Det er jo et Ordspil, en spøgefuld Spillen
med et dobbelttydigt Ord. Først tages Ordet „existere“ i rent abstract
Betydning som det at være til, og saa hedder det med Rette, at en Bevidsthed
ikke kan være theoretisk uden at existere, derpaa tages det i den særegne,
nylig indførte Bemærkelse af Livet for de praktiske Ideer og nu hedder det om
Existensen, at den selvfølgelig maa sættes øverst. Ak! det „øverst og nederst“
det er Prof.\ Nielsen en dyrebar Betegnelse; kunde de hinanden modsigende
Ytringer virkelig forsones, saa var det saa vist eet og det samme, hvor paa
Lav jeg satte dem. Hvis Troens Superioritet ikke maa yttre sig i en Omformning
af Viden, er det ikke godt at fatte, hvori den bestaaer. Men altsaa: fordi man
ei kan theoretisere uden at existere ɔ: spise, drikke, sove etc., derfor maa
% `ɔ:' --- the old Danish id-est mark (reversed c + colon, U+0254), read at
% 900 dpi: the glyph is an unmistakable reversed c, not an o and not a plain
% colon.  Google drops the reversed c and leaves a bare `:'; KB reads `o:'.
% The preamble already maps U+0254 to \reflectbox{c}; the same mark occurs in
% the pp. 25--35 fragment.
en positiv Troes Paradoxer sættes over Menneskeaandens fornuftige Resultater.
Men ligemeget med den svage Begrundelse. Lad os blot undersøge i og for sig,
om det Theoretiske ogsaa lader sig opgive! Ja, hvis Theorien virkelig var eet
med den systematiserede Videnskab og denne igjen eet med saa og saa mange
trykte i Papir- eller Læderbind indbundne Blade, da lod det sig gjøre at lade
Videnskaben ligge og lade Theorien være Theorie. Men hvis det ikke forholder
sig saaledes, hvis der i den simpleste Samtale mellem en Bonde og en Fisker,
hvis der i den jævneste Udtalelse af to enfoldige Tanker, mellem hvilke der er
Sammenhæng og i hvis Række%
% --- p. 69 ---
% The page turn falls INSIDE the word `Række-/følge' (p. 68 ends `Række-',
% p. 69 opens `følge'), so the marker is set mid-word: the `%' after `Række'
% eats the newline, and the comment lines contribute nothing, so `Række' joins
% directly to `følge' across the page turn.
% No footnote.  No letterspacing: spacing.py flagged `single? Vil-' (1.21),
% which is the hyphenation fragment of „Vil-/liens Concentration' --- ordinary
% spacing at the image.
\opage{69}følge der er Consequens --- hvis der, siger jeg, i disse ganske
ligefremme og dagligdags Phænomener er en spirende Videnskab, en elementær
Theorie, hvad da? Nei, Videnskabens Væsen er ikke afhængig hverken af Papir og
Trykkersværte eller af System og Fuldstændighed, Theorien er tilstede overalt,
hvor en fornuftig Erkjendelse finder Sted, den være nok saa abrupt, nok saa
enfoldig. At opgive Theorien gaaer da ikke saa let.

En anden Form for den samme Tankegang forekommer hos Dr.\ Heegaard anf.\ Sted
p.\ 447. Fremstillingen er der lagt an paa at charaktisere den „Villiens
% PRINTER'S ERROR, transcribed as printed: `charaktisere', not
% `charakterisere'.  Read at 900 dpi --- c-h-a-r-a-k-t-i-s-e-r-e, no `er'
% syllable.  Google reads `charaktisere' too; KB reads `charaklisere'.
Concentration ud over det Rationelle“ der findes hos den Paradoxet Troende.
Den er saa stærk, at under den alle Hensyn til Theorien forsvinde, „ja,“ siges
der, „end ikke den rationelle Ethik paa sit Maximum kan her komme i
Betragtning“. Godt! saalænge Individet ikke har denne Concentration, er det
kun halvveis troende, og en halv Tro er ingen Tro; saasnart det derimod har
hævet sig til Troens Høide, har det jo tabt al Theorie aldeles afsyne. Hvilken
tydeligere Udtalelse af Dualismens Umulighed kan man vel ønske end denne, der
er givet af dens egen Begrunder og sat i Pennen af en af dens blindeste
Forsvarere!

Til denne Vending slutter sig saa den, hos Kierkegaard imponerende, men her
uheldige, at den sande Troende ingen Tid har til at beskjæftige sig med
% --- p. 70 ---
% ONE footnote, mark `*)' --- SINGLE asterisk, read at 900 dpi both at the
% mark (`Theorie*).') and at the foot.  The Google layer renders the body mark
% as `** )'; that is OCR noise, there is no second note on this page.
% CAPITAL Æ: `Æren' read at the image, the ligature unambiguous.  Google reads
% `Eren', exactly as the audit predicted (ORTHOGRAPHY.md §3a); KB has `Æ ren'.
% `Deel' here, against `Del' on pp. 66 and 74 --- as printed, not normalised.
% `ligervis' and `naturligvis' both print a SINGLE i (ORTHOGRAPHY.md §1).
% No letterspacing: spacing.py's two flags (`an-', `op-') are hyphenation
% fragments; the footnote was read for letterspacing BY EYE and is set in
% ordinary spacing throughout.
\opage{70}den videnskabelige Theorie\fnmark{*)}{See Heegaard p.\ 395, smlgn.\
% `See' with double e in this note, against `Se' in the note to p. 72 below.
% Both read at the image; the book is inconsistent and is not normalised.
% `p. 395' --- read at 900 dpi: the first digit is the flat-topped old-style 3,
% distinct from the 5 that closes the same number.  KB reads `595'.
Nielsen: Evangelietroen og Theologien p.\ 155.}. Hvis denne Vending ikke
opklarer Andet, saa viser den idetmindste, at vore Theoretici, ligervis som de
gamle Vikingekonger, der tilfredse med Æren gjerne afstode fra deres Deel af
Byttet, høimodigt for deres egne Personer gjøre Afkald paa den Tro, de som
Feltherrer saa tappert tilkjæmpe de Andre.

Endnu en Udvei erindrer jeg at have hørt Prof.\ Nielsen antyde paa sine
Forelæsninger; da jeg ikke har noget trykt Document at støtte mig til, berører
jeg den naturligvis med et vist Forbehold, om jeg end ikke tvivler paa, at
Prof.\ N. vil vedkjende sig at have fulgt den Tankegang, hvorpaa det her
kommer an. Prof.\ omtalte det Bibelsted, hvor der fortælles, at Himmelen
aabnede sig, og Aanden foer ned som en Due, det samme Sted, som Hr.\ A.\ C.\
Larsen har anført i sin Bog „Samvittighed og Videnskab“ pag 59. Det
% `pag 59' --- NO stop after `pag' here, against `pag. 95' in the note to
% p. 72.  Read at the image; transcribed as printed.
Spørgsmaal blev opkastet, om Videnskaben nu her turde anbringe sin profane
Kritik, og kræve den Troende til Regnskab for, hvorledes han bærer sig ad med
at fastholde en Hvælving, naar han veed, der ingen er, med at troe, den
aabnede sig, med at opfatte Duen som Andet og Mere end et rent Symbol, naar
han veed, at Aanden umuligt kryber sammen til en Due. Dette var Spørgsmaalet,
om end udtrykt i andre Ord. Prof.\ dadlede nu den Aandløshed, der
% --- p. 71 ---
% No footnote, no letterspacing.
\opage{71}hang sig i det enkelte Billede, der ikke forstod, at for Troen vare
disse Billeder i Flyden og Strømning, i en uophørlig Bevægelse, under hvilken
det var ligesaa umuligt for Kritiken at gribe som for Forstanden at fatte dem.
Naar man hører Sligt første Gang, da skammer man sig over at have været saa
aandløs, man føler sig tilmode som en beskjæmmet Chicaneur; men ved nøiere
Eftertanke viser det sig ogsaa her, at Standpunktet kun lader sig hævde, naar
% `her,' --- COMMA.  KB reads a full stop; the image has a comma.
dets Eiendommelighed opgives, og den stærke Spænding mellem Troens og Videns
Synsmaader slappes. Nylig saa vi, at man hjalp sig ved at kaste Theorien
tilside, her synes det omvendt at gaae ud over Troen. Thi den klare,
historiske Virkelighed, som den Troende maa tillægge de hellige
Kjendsgjerninger, hvis han ikke til sin Forundring vil see sig enig med
Rationalisten, synes her at vige Pladsen for en mystisk Halvvirkelighed, en
ideel Realitet, der Intet har at sige der, hvor det netop gjælder om Existens
og Factum.

Men tilbage til Prof.\ Nielsens Opofrelse eller mulige Opofrelse af Theorien
til Bedste for det praktiske Liv. Hvad er det vel, der skal opnaaes derved?
ganske for Intet pleier man ikke at gjøre endog blot et muligt Offer.

Det er Villiens Oprindelighed og Selvtilstrækkelighed, der skal sikres,
forsaavidt er Offeret forklarligt nok. Uforklarligt er det derimod, at noget
Synderligt kan antages at være vundet derved. Thi det maa dog vel staae fast,
at Villien i og for sig er en
% --- p. 72 ---
% ONE footnote, mark `*)' (single asterisk, read at the image at the mark and
% at the foot).  The note was read for letterspacing BY EYE: ordinary spacing.
% No letterspacing in the body either.
\opage{72}ganske ligesaa immanent, ligesaa lidet overnaturlig Ytring af det
Menneskelige som Tanken. Det er jo dog en Opfattelse, der tilhører en forældet
Tids Naivetet, at visse af Menneskelivets og visse af Jordelivets Phænomener
eller Ytringer ere mere supranaturale end de andre. Formedelst en saadan
Opfattelse rystede mangen gammel troende Jøde paa Hovedet, da Olien i den
kjøbenhavnske Synagoges evige Lampe blev afløst af Gas. Thi Gas, det er et
revolutionært Element, høist naturligt, høist aandigt, og det Lys, som den
spreder, er stærkt og skarpt, Olie derimod er en blid tyktflydende Vædske, der
har gammel Hævd paa sig som den eneste rette Lysbringer og som desuden er i
Familie med Salvelse og alt andet Deslige. Skal nu Tanken paa samme Maade
bestemmes som profan, Villien derimod udgives for at være i Slægt med det
Positive? Er Villien da Tro eller Villien det Organ, hvormed Troen gribes? Ja
det paastaaes, men denne Paastand er aldeles falsk. Villien er ikke saadan
halvt om halvt en transscendent og aabenbaret Magt, den er rent menneskelig,
rent fornuftig. Er Troens Gjenstand et Under, saa er Troen selv lige saa fuldt
et Under\fnmark{*)}{Se „Philosophiske Smuler“ pag.\ 95.}, og Villien er ikke
% `Se' (single e) and `pag.' (with stop) here, against `See' and `pag' on
% p. 70 --- both as printed.
Underet nærmere end Tanken. Men paa dette Punkt gjør Prof.\ Nielsen sig
skyldig i en Række af Tilsnigelser og han, der saa haardnakket har bekjæmpet
den speculative Dogmatiks Sophismer, fremsætter en
% --- p. 73 ---
% CAPITAL Æ: `Ære' read at 900 dpi.  Google reads `Are'; KB reads `Æ re'.
% `Ringhed' --- the R's leg is broken in the inking and Google reads `Kinghed';
% at 900 dpi the bowl and the shoulder of an R are clear, and KB has `Ringhed'
% (ORTHOGRAPHY.md §3e).
% `for megen' is TWO words at the image; KB runs them together.
% No footnote.  No letterspacing: spacing.py's `single? ikke' (1.21) is the
% `maaskee ikke for megen' of this line, ordinary spacing at the image.
\opage{73}Sophisme, maaskee mindre blændende, men ikke mindre forkastelig end
dens. Det hedder i „Bemærkningerne“: „Skulde Nogen imidlertid være saa livlig
i Hovedet og saa let i Tankerne, at det virkelig løber rundt for ham, hvad der
skal staae øverst og hvad der skal staae nederst, Troen eller Viden, men af
lutter Livlighed forvexler Øverst med Nederst og Nederst med Øverst, lad ham
forsøge at simplificere Forholdene ved at kaste Troen, Villiesidealet,
Charakterprincipet bort og leve paa den rene Theorie: han undgaaer dog ikke
Fordoblingen“. Hvilken Fart! hvilke Spring som med Syvmilestøvler! fra
Existensen til Troen, fra Troen til Villiesidealet og Charakterprincipet! Jeg
arrogerer mig maaskee ikke for megen Ære ved at antage, at det er min Ringhed,
som er ment med den fortumlede Livlige, og at den anførte Passus er fremkaldt
ved en Ytring af mig i „Fædrelandet“ for 16de Februar 1866. Men hvis saa er,
hvor har jeg vel talt om at kaste først Troen i ren Almindelighed, saa
Villiesidealet og Charakterprincipet bort? Saa lidt som det er mig, der
anbefaler den Kur, at amputere Hovedet for at frelse Livet, saalidet ønsker
jeg nogen anden Lemlæstelse af det Menneskelige, hvorved intet Andet opnaaes
end at Feberen --- paa Holbergsk --- forlader den Syge. Hvad Andet har vel jeg
% Both dashes are em-dashes at the image; Google renders the first as `-' and
% the second as `--'.
i mine Bladartikler forsøgt at godtgjøre, end at to absolut uensartede
Principers Forening i samme Bevidsthed er umulig? Hvor staaer det skrevet, at
den, der opgiver den positive Religion,
% --- p. 74 ---
% No footnote, no letterspacing.
% `Del' (not `Deel') --- read at the image; cf. p. 70.
\opage{74}opløser det Religiøse selv i det Videnskabelige? Een saadan
Personlighed som den amerikanske Præst Theodor Parker er Vidnesbyrd nok om
denne Paastands Løshed. Og hvordan vil man vel overbevise om, at den, der
bryder med Orthodoxien, svigter Villiesidealet, at den, der vil bevare
Charakterens Eenhed i det theoretiske og i det praktiske Liv, at han
bortkaster Charakterprincipet? Hvordan vil man godtgjøre, at det er paa denne
% `bortkaster Charakterprincipet' --- no stop after `bortkaster'.  KB inserts
% one; the image has none.
Side Charakterløsheden, paa hin derimod, at Fastheden findes? Den nye Lære maa
her, som saa tidt, henvise til Livet og hvad Livet lærer. Hvad lærer da Livet,
hvad Svar giver det? Ja her, som bestandig, er Vanskeligheden den, at Livet
lærer Enhver et Forskjelligt. Men kan der fra den nye Læres Side vel indvendes
Noget, naar En for sin Del erklærer, at de sandeste Charakterer, den reneste
Villie har han fundet paa den Side, hvor efter Theorien de „intellectuelle
Amphibier“ skulde opholde sig, men derimod de strømmende, de i intellectuel
Begeistring vibrerende Aander der, hvor efter Theorien Charakterens Fasthed
netop skulde findes. Af egen Drift vilde vel Ingen bruge et saa svagt, saa
fruentimmeragtigt Argument, men jeg seer ikke rettere, end at Prof.\ Nielsen
selv er den, som har hjemlet Brugen.

Der forekommer i Aandens Verden hyppigt Phænomener af en saa afgjørende
Betydning, at de have Krav paa vor Interesse, ligegyldigt om de opdukke i nok
saa stor en Afstand fra os eller midt i vort eget
% --- p. 75 ---
% No footnote.  No letterspacing: spacing.py's `single? Man' (1.17) is the
% sentence-opening `Man har beraabt sig', ordinary spacing at the image with a
% wide justification gap before it.
\opage{75}Fædreland; der gives andre som vi ingen Opmærksomhed vilde skjænke,
hvis de vare opstaaede i China eller i Amerika eller i Tydskland, men som vi
føle os forpligtede til at anerkjende eller bestride, ene og alene fordi deres
Nærhed i Tid og Rum udruster dem med en Magt, som deres indre Værd aldrig i og
for sig vilde forlene dem. Prof.\ Nielsens Lære om Forholdet mellem Religion
og Videnskab hører til den sidste Klasse. Grunden til, at vi Danske bør agte
paa den, er den, at den er dansk. Dog, hvad siger jeg, dansk! er denne
unaturlige, med sig selv splidagtige Lære dansk? Den er det kun i Ordets
bogstavelige Forstand; i dybere og sandere Mening dansk er den rigtignok meget
langt fra at kunne kaldes, saalænge man hævder den danske Tænkning Sundhedens
og Klarhedens Egenskaber, saalænge man er enig om, at den danske Tankegang vel
ikke er mørkeræd som saa tidt den franske, men endnu mindre dunkel som saa
hyppigt den tydske og allermindst er lyssky eller lovløs. Den er da blot dansk
i Henseende til sin Oprindelse og forsaavidt som dens Tilblivelse kun er
begribelig ud fra givne historiske Forudsætninger paa dansk Grund. Man har
beraabt sig paa dette Forhold til Fortiden for at hævde den nye Læres
Nødvendighed, for at fremstille enhver Opposition imod den som et afmægtigt
Forsøg paa at holde Tidens eget vældige Hjul tilbage, men man har,
beskjæftiget, som man var, med at føle Verdenshistoriens Puls, glemt at,
% PRINTER'S ERROR, transcribed as printed: a comma after `glemt at'.  Read at
% the image; both layers have it too.  `glemt at, Spørgsmaalet her ikke er,'
Spørgsmaalet her ikke er,
% --- p. 76 ---
% LAST PAGE OF THE BOOK.  Structure at the image, top to bottom:
%   1. the chapter's last paragraph, running on from p. 75 and ending
%      `sin Arveret til Sandhedens Trone.';
%   2. a short CENTRED RULE, about one-sixth of the measure, set in the white
%      space below it --- reproduced here as \rule;
%   3. the unheaded closing note, set as an ORDINARY INDENTED PARAGRAPH in the
%      body type (no heading, no smaller face, no display indentation), opening
%      with the LETTERSPACED lead-in `Efterskrift:' followed by an em-quad of
%      space and then `En, efter at dette Skrift var'.  The letterspacing was
%      confirmed at 900 dpi: E-f-t-e-r-s-k-r-i-f-t carries the same widened
%      gaps as the book's other Sperrsatz, and the colon sits after the
%      trailing letterspace.  It is the ONLY emphasis in pp. 66--76.
%      spacing.py did NOT detect it --- another instance of the fit failing on
%      a short isolated run.  The colon is set inside \emph{} because the gap
%      before it is the letterspacing's own trailing space.
%   4. a DOUBLE RULE (two parallel lines, about one-third of the measure,
%      centred) closing the page well below the note.  Not reproduced: it is
%      the printer's terminal ornament, not part of the text.
%   5. `26 JY 67' --- British Museum accession datestamp, NOT text.
%   6. `6' at the foot --- the printer's signature, NOT text.
% `Pièce' carries a GRAVE accent: read at 900 dpi, the stroke descends left to
% right.  The KB layer reads an acute (`Piéce'); Google and the image have the
% grave.  It is set in the ordinary roman with no contrasting face, so it takes
% no markup --- cf. `trinité chrétienne' p. 56 and `n'observe' p. 65.
% The quotation marks around the title are „…“ at the image; Google collapses
% the opening one to `,,' and the closing one to `"'.
% No footnote on this page.
\opage{76}om Theorien er fattelig som et naturligt Foster af en tidligere
Tilstand eller forklarlig som Undflyen af en dobbelt Ild: men om den er
fyldestgjørende, om den er sand, om den Bram, hvormed den af Tilhængere
forkyndes, paa nogen Maade er berettiget, og man har ikke betænkt, at enhver
eklektisk Lære med samme Lethed kan godtgjøre sin ægte Herkomst og sin
Arveret til Sandhedens Trone.

\begin{center}
\rule{0.16\textwidth}{0.4pt}
\end{center}

\emph{Efterskrift:} En, efter at dette Skrift var givet i Trykken, udkommen
Pièce af Prof.\ R.\ Nielsen „Svar til Hr.\ Dr.\ phil.\ Pastor Zeuthen“
foranlediger mig ikke til nogen Tilføielse eller Rettelse i det Skrevne. Den
indeholder for største Delen kun Optryk af Forf.s Udfald mod Martensens
Dogmatik.

\end{document}
